% !TeX program = xelatex
\documentclass[12pt]{nextbook}

\UseTemplateSet{
  globalfonts = {libertinus},
  features = {hyperlinks}
}

\usepackage{paracol}
\usepackage{xcolor}
\usepackage{eso-pic}
\usepackage{needspace}

\vbadness=10000
\hbadness=10000
\hfuzz=3pt
\emergencystretch=1.5em

\makeatletter
\AtBeginDocument{%
  \@ifpackageloaded{hyperref}{%
    \pdfstringdefDisableCommands{%
      \def\ruby#1#2{#1}%
      \def\ReaderVocabText#1{#1}%
      \def\ReaderVocabTerm#1{#1}%
    }%
  }{}%
}
\makeatother

\providecommand{\ReaderColumnSep}{6mm}
\providecommand{\ReaderColumnRatio}{0.75}
\providecommand{\ReaderVocabBaselineskip}{20pt}
\providecommand{\ReaderVocabText}[1]{\textbf{#1}}
\providecommand{\ReaderVocabTerm}[1]{\textbf{#1}}
\providecommand{\ReaderVocabSeparator}{:}
\providecommand{\ReaderVocabItem}[2]{\noindent\ReaderVocabTerm{#1}\ReaderVocabSeparator\ #2\par}

\setlength{\columnsep}{\ReaderColumnSep}
\setlength{\columnseprule}{0pt}
\columnratio{\ReaderColumnRatio}
\swapcolumninevenpages

\newcounter{readersection}
\newwrite\vocabfile
\newif\iffirstreaderlabelwritten
\newlength{\readercolumnswidth}
\newlength{\readertitlegutter}
\setlength{\readercolumnswidth}{\dimexpr\textwidth-\columnsep\relax}
\setlength{\readertitlegutter}{.75\columnsep}

\makeatletter
\newcommand{\readerseparatorline}{%
  \if@mainmatter
    \AtTextLowerLeft{%
      \ifodd\value{page}%
        \put(\LenToUnit{\dimexpr .75\readercolumnswidth+.5\columnsep\relax},0){%
          \color{black!35}\rule{0.3pt}{\textheight}%
        }%
      \else
        \put(\LenToUnit{\dimexpr .25\readercolumnswidth+.5\columnsep\relax},0){%
          \color{black!35}\rule{0.3pt}{\textheight}%
        }%
      \fi
    }%
  \fi
}
\makeatother

\AddToShipoutPictureBG{\readerseparatorline}

\newcommand{\currentvocabfile}{\jobname-vocab-\thereadersection.voc}
\newcommand{\beginvocabwrite}{\immediate\openout\vocabfile=\currentvocabfile}
\newcommand{\finishvocabwrite}{\immediate\closeout\vocabfile}
\newcommand{\vocabitem}[2]{\ReaderVocabItem{#1}{#2}}
\newcommand{\printvocablist}{%
  \footnotesize
  \setlength{\parindent}{0pt}%
  \setlength{\parskip}{0pt}%
  \setlength{\baselineskip}{\ReaderVocabBaselineskip}%
  \raggedright
  \input{\currentvocabfile}%
}

\newcommand{\readersectiontitle}[1]{%
  \par
  \begingroup
  \ifodd\value{page}%
    \leftskip=\readertitlegutter%
    \rightskip=0.25\textwidth plus 1fil%
  \else
    \leftskip=\dimexpr .25\readercolumnswidth+.5\columnsep+\readertitlegutter\relax%
    \rightskip=0pt plus 1fil%
  \fi
  \section{#1}%
  \par
  \endgroup
}

\NewDocumentCommand{\voc}{omm}{%
  \ReaderVocabText{#2}%
  \IfNoValueTF{#1}{%
    \immediate\write\vocabfile{\string\vocabitem{#2}{#3}}%
  }{%
    \immediate\write\vocabfile{\string\vocabitem{#1}{#3}}%
  }%
}

\newenvironment{readersection}[1]{%
  \stepcounter{readersection}%
  \beginvocabwrite
  \Needspace{10\baselineskip}%
  \iffirstreaderlabelwritten\else
    \label{reader:firstpage}%
    \global\firstreaderlabelwrittentrue
  \fi
  \readersectiontitle{#1}%
  \begin{paracol}{2}
}{%
  \finishvocabwrite
  \switchcolumn
  \printvocablist
  \end{paracol}
}


\title{Liber elementarius linguae Latinae}
\author{Kelvin Yang}
\date{\today}

\begin{document}
\frontmatter
\maketitle
\tableofcontents
\mainmatter
\chapter{Textus elementarii}
\begin{readersection}{Disce et Vīve}
\voc[discō]{Discō}{to learn} et \voc[audiō]{audiō}{to hear}. \voc[bene]{Bene}{well} audiō, bene discō. \voc[legō]{Legō}{to read} et \voc[scrībō]{scrībō}{to write}; bene legō, \voc{sed}{but} \voc[male]{male}{badly} scrībō. \voc[nōn]{Nōn}{not} male scrībere \voc[dēbeō]{dēbeō}{must, should}; bene scrībere dēbeō.

\voc[dīco]{Dīcō}{to say}: “Lege et scrībe; disce et audī.” Dīcitis: “Discimus, audīmus, legimus, scrībimus.” \voc[semper]{Semper}{always} discere dēbēmus.

\voc[videō]{Videō}{to see, look} et \voc[sentiō]{sentiō}{to feel, think}. \voc{sī}{if} bene audīs, bene sentīs; sī male audīs, male sentīs. \voc[amō]{Amō}{to love} discere, audīre, legere et scrībere.

\voc[dō]{Dō}{to give} et \voc[capiō]{capiō}{to take}: das et capis; damus et capimus. \voc[vīvō]{Vīvō}{to live}\voc{-que}{and} discō. Vīve bene; discite bene; legite bene; scrībite bene.
\end{readersection}

\begin{readersection}{Puella et Magistra}
\voc[puella]{Puella}{girl} \voc[ad]{ad}{to, toward} \voc[magistra]{magistram}{female teacher} \voc[veniō]{venit}{to come}. Magistra puellam videt et audit. Puella legit et scrībit; magistra dīcit: “Bene legis, sed semper discere dēbēs.”

Magistra \voc[historia]{historiam}{history} legit. Historia est magistra \voc[vīta]{vītae}{life}. Puella historiam audit et bene sentit. Magistra puellae dīcit: “Vīta nōn semper facilis est; sed sī bene discis, bene vīvis.”

In \voc[Rōma]{Rōmā}{Rome} sunt \voc[incola]{incolae}{inhabitant}. Incolae Rōmae \voc[patria]{patriam}{fatherland} amant. \voc[poēta]{Poēta}{poet} dē patriā scrībit, et puella verba poētae legit. Poēta \voc[victōria]{victōriam}{victory} patriae narrat.

\voc[agricola]{Agricola}{farmer} in \voc[terra]{terrā}{earth, land} est. Agricola terram amat et \voc[aqua]{aquam}{water} habet. \voc[nauta]{Nauta}{sailor} \voc[ē]{ē}{out from} aquā venit et ad terram venit. Agricola \voc[cum]{cum}{with} nautā dē aquā et terrā saepe dīcit.

Puella \voc[sine]{sine}{without} magistrā male legit; cum magistrā bene discit. Magistra puellae aquam dat et dīcit: “In terrā et in aquā est vīta. \voc[etiam]{Etiam}{and also, even} agricolae et nautae vītam amant.” Puella audit, videt, sentit.

Magistra \voc[ā]{ā}{away from} puellā librum capit et historiam scrībit. Puella dē Rōmā, dē patriā, dē terrā, dē aquā legit. In historiā est vīta; in vītā est historia. Puella \voc[habeō]{habet}{to have} historiam et magistram; puella saepe legit, saepe scrībit, saepe discit.

\end{readersection}

\begin{readersection}{Liber et Oppidum}
\voc[magister]{Magister}{teacher} \voc[liber]{librum}{book} habet. Liber est \voc[bonus]{bonus}{good} et \voc[Rōmānus]{Rōmānus}{Roman}. Magister librum legit, puerī et puellae audiunt. Magistra quoque audit et scrībit.

In librō est \voc[oppidum]{oppidum}{town}. Oppidum est \voc[magnus]{magnum}{big, great}; in oppidō est \voc[templum]{templum}{temple}. In librō etiam est \voc[bellum]{bellum}{war}, sed magister dīcit: “Bellum malum est.” \voc[populus]{Populus}{people, nation} Rōmānus templum amat. In templō \voc[deus]{deus}{god} est, et \voc[dea]{dea}{goddess} quoque est. Populus deum et deam \voc[laudō]{laudat}{to praise}.

\voc[puer]{Puer}{boy} librum capit et legit: “\voc[ante]{Ante}{before, in front of} templum est populus; \voc[apud]{apud}{at, among} populum est poēta; \voc[post]{post}{after, behind} templum est aqua; \voc[per]{per}{through, via} oppidum venit \voc[vir]{vir}{man}.” Magister puerum laudat, quia bene legit.

\voc[fēmina]{Fēmina}{woman} ad templum venit. Vir cum fēminā ambulat, puer cum magistrō legit. Magister est vir doctus; magistra est fēmina docta. Puer magistrum audit et dīcit: “Liber est pulcher.” Tum magister respondet: “Ita est; liber \voc[pulcher]{pulcher}{beautiful} est, sed liber quoque difficilis est.”

In librō sunt \voc[multus]{multī}{many} puerī, multae puellae, multī virī, multae fēminae. Populus Rōmānus magnus est. Historia Rōmāna magna est. Historia \voc[Graecus]{Graeca}{Greek} quoque in librō est. Puer historiam Graecam legit, sed hodiē historiam Rōmānam scrībit.

Magister dīcit: “\voc[amīcus]{Amīcus}{friendly} puer librum dat; malus puer librum nōn dat.” Puer bonus librum puellae dat. Puella librum capit et legit. Magister puerum laudat: “Puer bone, bene vīvis.” Puer respondet: “Magister bone, discere amō.” Magister rīdet et dīcit: “Qui bene discit, \voc[beātus]{beātus}{happy, fortunate, blessed} est; qui male discit, nōn beātus est.” \voc[malus]{Malus}{bad} liber male docet, sed liber bonus bene docet.

\end{readersection}

\begin{readersection}{Ego et Tū}
\voc[hodiē]{Hodiē}{today} ego per \voc[via]{viam}{road} veniō. Tū quoque venīs. Magister \voc[adsum]{adest}{to be present, assist}, magistra adest, puerī et puellae adsunt. Ego bene \voc[valeō]{valeō}{to be strong, be well}; tū nōn male valēs. Magister nōs videt et dīcit: “Vos bene discitis.”

Magister mihi \voc[exemplum]{exemplum}{example} dat. In exemplō est \voc[verbum]{verbum}{word}. Ego verbum legō et scrībō; tū verbum audīs et dīcis. Magister est \voc[doctus]{doctus}{learned}. Is \voc[scientia]{scientiam}{knowledge, science} amat, et scientia mihi \voc[prōsum]{prōdest}{to be useful}.

Ego magistrō \voc[crēdō]{crēdō}{to trust, believe}, \voc[quia]{quia}{because} magister bonus est. Tū quoque magistrō crēdis. \voc[quoque]{Quoque}{also, too} puella magistrō crēdit. Sī bene audīmus, discere \voc[possum]{possumus}{to be able}. Sed sī liber mihi \voc[dēsum]{dēest}{to lack}, male legō; sī liber tibi dēest, male scrībis. Nōn bene discimus \voc[nisi]{nisi}{unless} librōs habēmus.

\voc[līberī]{Līberī}{children} puerī puellaeque sunt. Līberī librōs capiunt, verba legunt, scientiam amant. Magister dīcit: “Qui scientiam amat, \voc[līber]{līber}{free} est.” Ego mē \voc[nōscō]{nōscō}{to come to know, learn}; tū tē nōscis; nōs nōs nōscimus; vos vōs nōscitis.

Sed puer \voc[labōrō]{labōrat}{to labor, be in distress}, quia verba male legit. Tum \voc[medicus]{medicus}{doctor} venit. Medicus puerum videt et dīcit: “Puer nōn male valet.” Puer respondet: “Valeō, sed verba difficilia sunt.” Magister verba legit \voc[sīcut]{sīcut}{just as} medicus puerum videt: bene et saepe. Ita puer verba difficilia \voc[vincō]{vincit}{to win}.

\voc[fortūna]{Fortūna}{fortune, luck} nōn semper adest, sed scientia saepe adest, sī semper discimus. \voc[crās]{Crās}{tomorrow} ego et tū iterum venīmus; nōs librōs legimus, vos verba scrībitis. \voc[sīc]{Sīc}{thus} discimus, sīc scrībimus, sīc vīvimus.

\end{readersection}

\begin{readersection}{Fābula in Scholā}
Magister in scholā librum habet. Puerī et puellae adsunt; etiam \voc[discipulus]{discipulī}{student} novī veniunt. Magister dīcit: “Hodiē \voc[fābula]{fābula}{story, tale} Rōmāna legitur et nārrātur. Audīte bene.”

In fābulā est \voc[ager]{ager}{field, farm}. In agrō sunt agricolae et paucī puerī. Prope agrum est \voc[mūrus]{mūrus}{wall, city wall}; post mūrum sunt \voc[castra]{castra}{encampment, camp}. Magister tabulam tenet et discipulīs agrum, mūrum, castra \voc[mōnstrō]{mōnstrat}{to show, point out}. Deinde discipulī rogant: “\voc[ubi]{Ubi}{where} sunt castra? \voc[cūr]{Cūr}{why} mūrus ante castra est?”

Magister respondet: “Castra ā mūrō servantur, quia ibi est \voc[perīculum]{perīculum}{risk, danger}. Sed agricolae nōn semper \voc[timeō]{timent}{to fear}. Agricolae agrum amant et in agrō labōrant.” Tum fābula ā magistrō \voc[narrō]{nārrātur}{to tell}. Discipulī audiunt, et verba magistrī bene intellegere cupiunt.

Magister discipulōs \voc[doceō]{docet}{to teach}; discipulī ā magistrō docentur. Liber ā puerō legitur; verba ā puellā scrībuntur; historia ā magistrā explicātur. Puer dīcit: “Magister, \voc[quōmodo]{quōmodo}{how} mūrus servātur?” Magister respondet: “Mūrus ab incolīs servātur; castra autem ā virīs Rōmānīs servantur.”

Tum alius discipulus interrogat: “\voc[quandō]{Quandō}{when} poena datur?” Magister dīcit: “Poena datur, sī populus male vīvit. Sed poena nōn semper est bona. \voc[poena]{Poena}{penalty, punishment} interdum timētur, interdum intellegitur.” Discipulī tacent, quia sententia difficilis est.

Magistra quoque adest et puerōs \voc[moneō]{monet}{to warn, advise}: “Nōlīte male legere; nōlīte male scrībere. Sī verba nōn intellegitis, \voc[quaerō]{quaerite}{to seek, ask}.” Tum discipulus quaerit: “Magistra, \voc[scīō]{scīs}{to know, understand}ne quid in fābulā agitur?” Magistra respondet: “Sciō; sed vos quoque scīre dēbētis.”

In tabulā duo verba scrībuntur: \voc[aequus]{aequus}{equal} et \voc[paucus]{paucī}{few, little}. Magister dīcit: “Paucī discipulī male audiunt, sed multī bene discunt. Labor discipulī et labor magistrī aequus esse nōn potest; magister docet, discipulus docētur.” Puer rīdet, sed magister eum iterum monet.

Deinde magister dīcit: “Nunc legite: castra \voc[contrā]{contrā}{against} oppidum sunt. Agricolae \voc[aut]{aut}{or} in agrō manent aut ad oppidum veniunt. Hīc liber est, \voc[hīc]{hīc}{here} fābula legitur, hīc verba intelleguntur.” Discipulī librum vident et fābulam legunt.

Post fābulam magister rogat: “Quid intellegitis?” Puella respondet: “Intellegimus perīculum, poenam, agrum, mūrum et castra.” Puer addit: “Intellegimus quoque quōmodo fābula nārrētur et cūr discipulī doceantur.” Magister laetus est et dīcit: “Bene. Qui discere amat, \voc[intellegō]{intellegit}{to understand, comprehend}; qui numquam audit, \voc[numquam]{numquam}{never} bene discit.”

\end{readersection}

\begin{readersection}{Hoc Cōnsilium}
\voc[nunc]{Nunc}{now} magister in scholā est. Hic magister librum habet; ille puer legit; ista puella scrībit. Magister dīcit: “Hic liber est bonus, ille liber est malus. Hoc verbum legite, illud verbum scrībite.”

Puer librum legit et \voc[cōnsilium]{cōnsilium}{plan, advice} magistrī audit. Magister dīcit: “\voc[ūnus]{Ūnus}{one} liber nōn \voc[satis]{satis}{enough} est; multī librī prōsunt. Sed \voc[nūllus]{nūllus}{no, none} liber prōdest, sī discipulus nōn legit.” Puella respondet: “Hoc cōnsilium mihi placet.”

Magister puerōs rogat: “\voc[uter]{Uter}{which of two} liber vōbīs placet, hic an ille?” Ūnus puer dīcit: “Hic mihi placet.” \voc[alter]{Alter}{the other of two} puer dīcit: “Ille mihi placet.” Magister tum rogat: “\voc[neuter]{Neuter}{neither} liber malus est?” Puerī respondent: “Neuter malus est; uterque bonus est.” Magister rīdet et dīcit: “Bene respondētis.”

Tum magistra intrat et \voc[servus]{servum}{slave} cum librīs videt. Servus librōs portat \voc[prō]{prō}{before, for, on behalf of} magistrō. Magister servō dīcit: “Librōs in mēnsā pōne.” Servus facit quod dominus iubet, quia \voc[dominus]{dominus}{master, lord} eius ibi stat. Sed magister puerīs dīcit: “Servus quoque homō est; nōn propter officium contemnendus est.”

In librō est fābula dē \voc[mundus]{mundō}{world, universe} et \voc[nātūra]{nātūrā}{nature}. Magister legit: “Mundus magnus est; nātūra multa facit.” Puella audit et dīcit: “Haec fābula mihi placet, quia dē vītā est.” Magister respondet: “Bene sentīs. \voc[sapientia]{Sapientia}{wisdom} est nātūram intellegere et vītam bene agere.” In hāc sententiā est \voc[animus]{animus}{soul, mind}: animus discipulī sapientiam quaerit.

\voc[alius]{Alius}{other} discipulus magistrum rogat: “Quid est sapientia?” Magister respondet: “Sapientia nōn est \voc[nimis]{nimis}{too much} dīcere; sapientia est satis intellegere. Qui nimis scrībit et nōn intellegit, male discit. Qui satis audit et bene legit, melius discit.” Hoc dictum discipulīs placet.

Deinde magister tabulārum exemplum mōnstrat. “\voc[tōtus]{Tōtus}{whole, entire} populus Rōmānus officium habet,” inquit. “Magister officium habet, discipulus officium habet, dominus officium habet, servus officium habet. \voc[officium]{Officium}{duty} est quod agere dēbēmus.” Puer dīcit: “Ego quoque officium habeō: discere dēbeō.”

Magister laudat puerum et dīcit: “Recte dīcis. \voc[agō]{Agis}{to do, make} quod discipulus bonus agit.” Puella quoque dīcit: “Ego verba scrībō; hoc officium meum est.” Magistra respondet: “Et ego vos doceō; hoc officium meum est.”

Tum dominus librum capit et interrogat: “\voc[ūllus]{Ūllus}{any} discipulusne hanc fābulam nōn intellegit?” Nūllus discipulus tacet; ūnus puer manum tollit et dīcit: “Ego nōn tōtam fābulam intellegō.” Magister respondet: “Bene facis quod rogas.” Deinde fābulam iterum legit.

Post lectionem magister dīcit: “Nunc vos ipsī fābulam \voc[faciō]{facite}{to make, do}. Ūnus dē mundō scrībat, alter dē nātūrā, alius dē sapientiā. Neuter puer male scrībat; uterque bene scrībat.” Discipulī verba audiunt et fābulam faciunt.

Tandem puella legit: “Hic mundus magnus est; haec nātūra pulchra est; hoc cōnsilium bonum est. Propter sapientiam bene vīvimus.” Magister dīcit: “Haec fābula mihi \voc[placeō]{placet}{to please, satisfy}.” Puella gaudet.

\voc[sōlus]{Sōlus}{only, alone} puer autem tacet. Magister eum videt et rogat: “Cūr tacēs?” Puer respondet: “Nōn taceō propter timōrem, sed propter cōnsilium. Primum audiō, deinde dīcō.” Magister laetus est: “Hoc quoque sapientia est.”

\voc[propter]{Propter}{near, because of} hoc cōnsilium discipulī minus errant. Hic legit, ille scrībit, iste audit; sed tōta schola discit. Magister tandem dīcit: “Nunc satis est. Valēte, puerī puellaeque; crās aliud exemplum legēmus.” Discipulī respondent: “Valē, magister bone.”

\end{readersection}

\begin{readersection}{Urbs, Rēx et Vēritās}
Magister hodiē librum novum habet. In librō est fābula dē \voc[urbs]{urbe}{city} Rōmānā. Discipulī audiunt, et magister dīcit: “Haec urbs magna est; in urbe sunt puerī, puellae, magistrī, nautae, agricolae et multī virī. Sed urbs sine lēgibus nōn bene vīvit.”

In fābulā \voc[rēx]{rēx}{king} in urbe habitat. Rēx populum videt et \voc[lēx]{lēgem}{law} scrībit. Lēx rēgis nōn semper bona est, sed rēx dīcit: “Lēx cīvitātī prōdest.” Tum magister discipulīs mōnstrat: “Vidēte: rēx, rēgis; lēx, lēgis. Haec sunt nōmina tertiae dēclīnātiōnis.”

Puer magistrum rogat: “Quid est \voc[cīvitās]{cīvitās}{city-state, state}?” Magister respondet: “Cīvitās est urbs et populus. In cīvitāte hominēs lēgēs habent.” Tum puella legit: “Cīvitās sine \voc[pāx]{pāce}{peace} labōrat; pāx cīvitātī prōdest.” Discipulī sententiam scrībunt.

Sub mūrō urbis stat \voc[homō]{homō}{human}. Is nōn est rēx, sed vir parvus. \voc[parvus]{Parvus}{little, small} homō magnum librum fert. Magister eum \voc[appellō]{appellat}{to call}: “Amīce, quid fers?” Homō respondet: “Librum et \voc[poēma]{poēma}{poem} ferō.” Tum puer dīcit: “\voc[ferō]{Ferre}{to bring, carry, bear} difficile est, sed librum ferre bonum est.”

Poēma dē \voc[mare]{marī}{sea} est. Nauta in marī vīvit et mare amat. Sed in poēmate etiam \voc[mors]{mors}{death} est, nam bellum urbi appropinquat. Puerī timent, sed magister dīcit: “Nōlīte timēre. Poēma nōn semper vērum est; interdum poēta dē morte scrībit, sed vītam laudat.”

In poēmate est quoque \voc[māter]{māter}{mother}. Māter fīlium amat et fīliō librum dat. Ea dīcit: “\voc[mēns]{Mēns}{mind} bona et \voc[virtūs]{virtūs}{virtue} magna sunt. Corpus sine mente nōn satis valet; mēns sine virtūte nōn beāta est.” Magister addit: “\voc[corpus]{Corpus}{body} vidētur; mēns autem intellegitur.”

Puella rogat: “Magister, quid est \voc[vēritas]{vēritās}{truth}?” Magister respondet: “Vēritās est quod nōn fallit. Rēx sine vēritāte malus est; cīvitās sine vēritāte labōrat; homō sine vēritāte nōn līber est.” Discipulī tacent, nam sententia gravis est.

Deinde magister dē \voc[mōs]{mōre}{custom, habit} Rōmānō dīcit. “Mōs populi saepe lēgibus \voc[differō]{differt}{to differ}. Lēx scrībitur; mōs in vītā hominum vidētur. Bonus mōs cīvitātī prōdest, malus mōs cīvitātem laedit.” Puer verba legit et intellegit.

Magistra intrat et \voc[opus]{opus}{work} discipulīs dat: “Scrībite dē urbe, dē rēge, dē lēge, dē pāce. Scrībite etiam dē homine parvō quī librum fert.” Discipulī opus accipiunt. Alius dē marī scrībit, alter dē morte, alia puella dē mātre et virtūte.

\voc[nam]{Nam}{because} tempus breve est, magister dīcit: “Nunc satis est. \voc[tempus]{Tempus}{time} fugit, sed ars manet.” Tum discipulī novum verbum audiunt: \voc[ars]{ars}{art}. Magister explicat: “Ars poētae est poēma facere; ars magistrī est docēre; ars discipulī est discere.” Puella respondet: “Ergō discere quoque opus est.”

Sub finem lēctiōnis liber \voc[sub]{sub}{under, below} mensā ponitur. Puer quaerit: “Cūr liber sub mēnsā est?” Magister rīdet et respondet: “Quia puerī nimis multa ferunt!” Discipulī rīdent. Sed magister iterum librum capit et dīcit: “Memoriā tenēte: animal, corporis, urbem, rēgī, pāce, lēgibus. Tertia dēclīnātiō difficilis est, sed vōbīs prōdest.”

In ultimā sententiā poēta scrībit: “Pāx cīvitātis, virtūs hominis, vēritās lēgis, ars poētae et amor mātris urbem servant.” Discipulī hanc sententiam legunt et sentiunt. Magister laetus est, quia discipulī nōn solum verba, sed etiam mentem poēmatis intellegunt.

\end{readersection}

\begin{readersection}{Cōnsul et Cīvēs}
In librō \voc[antīquus]{antīquō}{ancient, oldtime} magister fābulam dē Rōmā legēbat. Discipulī audiēbant et verba scrībēbant. Tum magister dīcēbat: “Ō discipulī, Rōma antīqua magna erat; ibi \voc[cīvis]{cīvēs}{citizen} multī vīvēbant, et \voc[iūs]{iūs}{right, law} cīvibus \voc[cārus]{cārum}{dear, beloved, costly} erat.”

In fābulā \voc[cōnsul]{cōnsul}{consul} in urbe stābat. Is \voc[iūstitia]{iūstitiam}{justice} amābat et \voc[lībertās]{lībertātem}{liberty} servābat. Cīvēs cōnsulem audiēbant, quia cōnsul dē iure et lībertāte bene dīcēbat. Tamen nōn omnēs cīvēs bene sentiēbant; paucī iūstitiam nōn amābant.

Sub \voc[caelum]{caelō}{heaven, sky} obscūrō \voc[hostis]{hostis}{enemy} ad urbem veniēbat. Hostis \voc[Gallus]{Gallus}{Gallic} erat, et cum hoste erant multae \voc[cōpia]{cōpiae}{plenty, troops}. Cōnsul cīvēs monēbat: “Hostēs prope sunt; sed lībertās nōbīs cāra est. Sī lībertātem amāmus, timēre nōn dēbēmus.”

Tum \voc[legiō]{legiō}{legion} Rōmāna ante mūrum \voc[stō]{stābat}{to stand}. Paucī puerī prope templum stābant et mīlitēs vidēbant. Magister discipulīs dīcēbat: “Legiō stābat, cīvēs manēbant, cōnsul imperābat, hostēs autem prope urbem veniēbant.”

Cōnsul nūntium ad legiōnem \voc[mittō]{mittēbat}{to send}. Nūntius per urbem currēbat et verba cōnsulis ferēbat. Cōnsul etiam cīvēs \voc[vocō]{vocābat}{to call}: “Cīvēs Rōmānī, hīc manēte; iūs et lībertātem servāte.” Cīvēs in urbe \voc[maneō]{manēbant}{to remain}; puellae et puerī cum magistrā apud templum erant.

Mox \voc[proelium]{proelium}{battle} incipiēbat. Hostēs pugnābant, et Rōmānī quoque \voc[pugnō]{pugnābant}{to fight}. Cōnsul nōn fugiēbat, sed ante legiōnem stābat. Is \voc[soleō]{solēbat}{used to} ita facere: ante perīculum manēbat, mīlitēs vocābat, et cīvēs monēbat. Ideō cīvēs eum amābant.

In mediō proeliō cōnsul signum \voc[inveniō]{inveniēbat}{to find}. Signum parvum erat, sed mīlitibus cārum erat. Cōnsul signum ad legiōnem mittēbat, et legiō animō fortī pugnābat. Tum cōpiae hostium iam recēdēbant, sed hostis Gallus tamen prope mūrum stābat.

Post proelium cōnsul in urbem redībat. Cīvēs eum laudābant, quia iūstitia et lībertās servābantur. Magister discipulīs dīcēbat: “Vidētisne? Cōnsul nōn sōlum bellum gerēbat; is etiam bonum exemplum \voc[creō]{creābat}{to create}. Nam qui iūstitiae servit, lībertātem servat.”

Puella sententiam scrībēbat: “Iūs cīvibus cārum erat; lībertās cīvitātī cāra erat; cōnsul iūstitiam amābat.” Puer autem dīcēbat: “Hostis pugnābat, tamen Rōmānī manēbant.” Magister respondēbat: “Bene scrībitis. Sīcut fābula docet, virtūs sine iūstitiā parum valet.”

Tum caelum clārum erat, et cīvēs in urbe vīvēbant. Cōnsul tamen semper vigilābat, quia hostēs interdum redībant. Discipulī fābulam legēbant et intellegēbant: in tempore antīquō cīvēs iūs, iūstitiam et lībertātem amābant; cōnsul cōpiās vocābat; legiō pugnābat; urbs manēbat.

\end{readersection}

\begin{readersection}{Iter Breve et Sōl Fēlīx}
Magister hodiē discipulīs fābulam legit. In fābulā est \voc[iter]{iter}{journey, trip} per agrōs et ad urbem. Iter nōn est \voc[longus]{longum}{long}, sed \voc[brevis]{breve}{short}. Puerī gaudent, quia iter breve facile amant.

In viā ambulant magister, discipulī et puer \voc[celer]{celer}{fast, speedy}. Puer celer ante omnēs currit, sed magister eum monet: “Nōlī nimis celer esse; iter sine cōnsiliō saepe \voc[difficilis]{difficile}{difficult} fit.” Puer audit et cum cēterīs manet.

Sub caelō est \voc[sōl]{sōl}{sun} clārus. Sōl super agrum lūcet, et via sub sōle vidētur. Magister dīcit: “Sōl omnibus hominibus \voc[commūnis]{commūnis}{common} est. Aqua quoque commūnis est, terra commūnis est, sed sapientia nōn semper commūnis est.”

Puellārum ūna librum \voc[teneō]{tenet}{to hold, possess}. In librō est \voc[nōmen]{nōmen}{name} urbis veteris. Puella legit: “Urbs habet nōmen antīquum; populus veterem memoriam tenet.” Magister addit: “Nōmen breve est, sed memoria urbis longa est.”

Mox discipulī ad mūrum veniunt. Illic sedet vir \voc[sapiēns]{sapiēns}{wise}. Is puerīs dīcit: “Homō \voc[mortālis]{mortālis}{mortal} est, sed virtūs et vēritās diu manent.” Puer hoc verbum audit et rogat: “Num omnis homō mortālis est?” Vir respondet: “Ita est; \voc[omnis]{omnis}{every, all} homō mortālis est.”

Magister discipulīs tertiae dēclīnātiōnis adiectīva mōnstrat: “Vir sapiēns, fēmina sapiēns, nōmen sapiēns nōn dīcimus; sed cōnsilium sapiēns dīcere possumus. Item: vir fortis, puella fortis, animal forte.” Tum in tabulā scrībit: \voc[fortis]{fortis}{strong, powerful}, forte.

Discipulus quaerit: “Estne haec rēs \voc[simplex]{simplex}{simple}?” Magister respondet: “Nōn semper. Forma brevis est, sed ratiō nōn est semper simplex. Tamen sī bene legitis, difficilis lēctiō vōbīs nōn nimis difficilis erit.” Discipulī rīdent, sed verba scrībunt.

Postea magistrum rogant: “Quis urbem \voc[regō]{regit}{to govern, rule}?” Magister respondet: “Rēx urbem regit, sed lēx quoque rēgem regit. Rēx fortis sine sapientiā nōn satis bonus est. Rēx sapiēns populum bene regit.” Discipulī intellegunt: potestas sine sapientiā perīculum est.

In portā urbis stat fēmina \voc[fēlīx]{fēlīx}{lucky, happy}. Ea puerīs aquam dat et dīcit: “Iter vestrum breve sit et fēlīx.” Magister respondet: “Grātiās tibi agimus; puerī enim aquam habēre dēbent.” \voc[enim]{Enim}{for, because} sōl calidus est, et discipulī aquam libenter capiunt.

Deinde omnēs \voc[super]{super}{above, upon} mūrum parvum spectant. Sōl super urbem est, nōmen urbis in portā scrībitur, et vir sapiēns iterum dīcit: “Vetus urbs nōn semper magna est, sed saepe sapientiam veterem tenet.” Magister discipulīs dīcit: “\voc[vetus]{Vetus}{old, ancient} est adiectīvum speciale: veteris, veterī, veterem, vetere.”

Sub finem itineris puer celer iam nōn currit. Is lentē ambulat et dīcit: “Iter breve fuit, sed verba multa erant.” Puella respondet: “Verba multa erant, sed lēctiō bona fuit.” Magister laetus est, quia discipulī nōn modo viam, sed etiam adiectīva tertiae dēclīnātiōnis didicērunt.

\end{readersection}

\begin{readersection}{Herī in Silvā}
\voc[herī]{Herī}{yesterday} \voc[pater]{pater}{father} \voc[fīlius]{fīlium}{son} ad \voc[silva]{silvam}{wood, forest} \voc[dūcō]{dūxit}{to lead}. Cum eō erat etiam \voc[uxor]{uxor}{wife}. Iter breve erat, sed puer gaudēbat, quia sōl super agrōs lūcēbat et silva pulchra erat.

In silvā pater fīliō veterem fābulam nārrāvit. “Rēx quondam hanc urbem servāvit,” \voc[inquit]{inquit}{he/she/it says or said}. “Hostēs urbem expugnāre cupīvērunt, sed cīvēs iūs et lībertātem amāvērunt.” Fīlius patrem \voc[interrogō]{interrogāvit}{to ask}: “Num hostēs urbem \voc[expugnō]{expugnāvērunt}{to conquer}?” Pater respondit: “Nōn; cōnsul legiōnem bene dūxit.”

Tum \voc[nōnnūllus]{nōnnūllī}{some, several} puerī in silvā apparuērunt. Eī librum portāvērunt et poēma recitāre \voc[coepī]{coepērunt}{to begin}. Poēma dē pāce erat, sed etiam dē bellō. Pater audiit et \voc[ait]{ait}{he/she/it says or said}: “Bene legitis; sed memoriā tenēte: bellum malum est.”

Sub arbore iacuit parvum animal. Puerī timuērunt, quia putāvērunt animal mortuum esse. Sed animal adhūc vīvēbat. Fīlius aquam dedit, et animal oculōs aperuit. “Nōn mortuum est,” fīlius inquit; “vix vīvit.” Pater animal tetigit et dīxit: “Nēmō hoc animal \voc[necō]{necāvit}{to kill}; fortasse perīculum effūgit.”

Puer fābulam audīvit et animum gravem habuit. “Bellum \voc[ōdī]{ōdī}{to hate},” inquit, “quia bellum hominēs et animalia necat.” Uxor eum laudāvit: “Bene sentīs. Qui bellum ōdit, pācem amat.” Pater quoque dīxit: “Ego multa bella in librīs lēgī, sed pācem semper amāvī.”

Mox omnēs ex silvā rediērunt. Fīlius \voc[meminī]{meminit}{to remember}: rēgem, cōnsulem, hostēs, animal parvum et verba patris. \voc[mox]{Mox}{soon, then} domum vēnērunt, et puer fōrmam animalis in tabulā scrībere cupīvit. \voc[fōrma]{Fōrma}{form, shape, beauty} animalis parva erat, sed pulchra.

\voc[vix]{Vix}{hardly, scarcely} domum intrāvērunt, cum fīlius librum aperuit et fābulam scrīpsit. “\voc[quamquam]{Quamquam}{although} silva obscūra erat,” scrīpsit, “ibi didicī pācem amāre et bellum ōdisse.” Pater legit et subrīsit. “Bene scrīpsistī,” ait. “Ūnus \voc[annus]{annus}{year} praeteribit, sed hunc diem semper memineris.”

\end{readersection}

\begin{readersection}{Domus et Rēs Pūblica}
\voc[diēs]{Diēs}{day} clārus erat, et magister discipulīs fābulam dē Rōmā nārrābat. In fābulā \voc[domus]{domus}{house, home, dwelling-place} antīqua apud viam erat. Ante domum stābat puer, librum tenēbat, et magistrum exspectābat. Magister vēnit et dīxit: “Hodiē dē quartā dēclīnātiōne discimus.”

In librō erat imāgō: vir manum tollēbat, et in \voc[manus]{manū}{hand} librum habēbat. Magister dīxit: “Vidēte: manus, manūs; manum, manū. Hic est ūsus quartae dēclīnātiōnis.” Puella manum suam mōnstrāvit et rīsit. Puer autem librum scrīpsit, quia verba memoria tenēre cupiēbat.

Tum magister \voc[exercitus]{exercitum}{army} Rōmānum in tabulā mōnstrāvit. \voc[magistrātus]{Magistrātus}{magistrate} exercitūī verba dabat, et exercitus ad urbem veniēbat. Magister dīxit: “Magistrātus est vir pūblicus; exercitus saepe magistrātibus pāret.” Discipulī audiēbant et castra, mūrōs, exercitum in tabulā vidēbant.

In eādem fābulā erat \voc[fidēs]{fidēs}{faith, belief}. Cīvēs magistrātuī crēdēbant, quia magistrātus iūstitiae serviēbat. “Sine fide,” inquit magister, “rēs pūblica labōrat. \voc[rēs]{Rēs}{thing} pūblica nōn est domus ūnius hominis, sed domus multōrum cīvium.” Puer hoc verbum scrīpsit et bene intellexit.

Puella interrogāvit: “Quid est \voc[ōrdō]{ōrdō}{order}?” Magister respondit: “Ōrdō est certus modus rērum. In domō est ōrdō; in exercitū est ōrdō; in linguā quoque est ōrdō.” Tum magister dē \voc[lingua]{linguā}{tongue, language} Latīnā dīxit: “Lingua Latīna multa verba habet, sed verba sine ōrdine male intelleguntur.”

Magister deinde dē \voc[ratiō]{ratiōne}{reason, rationality} dīxit. “Ratiō hominibus prōdest. Sī ratiō adest, animus melius iūdicat; sī ratiō dēest, timor saepe vincit.” Puer rogāvit: “Ergō ratiō est quasi lūx mentis?” Magister respondit: “Bene dīcis; ratiō mentem illūstrat.”

Sub arbore sedēbat senex. Is \voc[sēnsus]{sēnsūs}{sense, sensation} dēbilem habēbat, sed \voc[spīritus]{spīritus}{spirit} fortis erat. Senex puerīs dīxit: “Oculī vident, aurēs audiunt, sed spīritus intellegere dēbet. \voc[vīsus]{Vīsus}{vision} sine ratiōne nōn satis est.” Discipulī verba senis audīvērunt et tacitī erant.

Magister addidit: “Hominēs \voc[ūsus]{ūsū}{use, experience} discunt. Puer saepe legit; diū legit; tandem bene legit. Medicus quoque ūsū valet; nauta ūsū mare nōscit; agricola ūsū terram amat.” \voc[diū]{Diū}{for a long time} discipulī hanc sententiam scrībēbant, quoniam verba nova erant.

Tum puella librum aperuit. In librō erat \voc[speciēs]{speciēs}{view, kind, species} pulchra urbis: templum, mūrus, via, domus, populus. Magister dīxit: “Speciēs urbis pulchra est, sed nōn omnis speciēs vēritātem mōnstrat. Interdum rēs pulchra vidētur, sed mala est.” \voc[interdum]{Interdum}{sometimes} discipulī verba magistrī nōn statim intellegēbant.

Puer rogāvit: “Num est \voc[spēs]{spēs}{hope} in fābulā?” Magister respondit: “Est. Spēs cīvibus adest, quia exercitus urbem servāvit et magistrātus iūstitiae crēdidit. \voc[quoniam]{Quoniam}{because, since} fidēs et spēs adsunt, cīvitās nōn cadit.” Puer laetus erat, quia fābulam iam melius intellegēbat.

Ad finem lēctiōnis magister tabulārum ōrdinem mōnstrāvit: “Caput, capitis; domus, domūs; manus, manūs; rēs, reī; spēs, speī.” Tum puer \voc[caput]{caput}{head} inclināvit et dīxit: “Caput meum plēnum est verbīs.” Magister rīsit: “Nōlī timēre; ūsū et tempore verba manēbunt.”

Post lēctiōnem discipulī ad domum rediērunt. Ille puer librum in manū portābat et dīcēbat: “Hodiē multa didicī: domum, manum, exercitum, magistrātum, rērum ōrdinem, ratiōnem, fidem et spem.” Puella respondit: “Et ego didicī linguam Latīnam nōn esse rem parvam, sed rem pulchram et \voc[potēns]{potentem}{potent, mighty}.”

\end{readersection}

\begin{readersection}{Avis Reperta}
Magister hodiē discipulīs dīxit: “In hāc lēctiōne participia discimus. Vidēte: mīles currēns, avis fugiēns, urbs occupāta, dōnum relictum.” Discipulī audientēs verba in tabulā scrībēbant.

Tum magister fābulam novam nārrāvit. In fābulā \voc[dux]{dux}{leader, commander, general} cum \voc[mīles]{mīlitibus}{soldier} per silvam ambulābat. Sub arbore erat parva \voc[avis]{avis}{bird}; avis vulnerāta nōn volābat, sed in terrā iacēbat. Dux avem vīdit et \voc[statim]{statim}{at once, immediately} mīlitēs vocāvit.

Ūnus mīles ad avem \voc[currō]{cucurrit}{to run}. Avis autem mīlitem videntem timuit et \voc[fugiō]{fugere}{to flee, escape} temptāvit. Magister dīxit: “Mīles currēns avem videt; avis fugiēns auxilium quaerit.” Puella verba scrīpsit: currēns, fugiēns.

Prope viam stābat \voc[lēgātus]{lēgātus}{legate, envoy}. Is dūcī dīxit: “Hostēs hunc \voc[locus]{locum}{place} paene \voc[occupō]{occupāvērunt}{to occupy}. Nōlīte hīc manēre.” Sed dux respondit: “Avem \voc[vīvus]{vīvam}{alive, living} relinquere nōn possumus.” Tum mīlitēs audientēs tacitī manēbant.

Dux mīlitēs \voc[iubeō]{iussit}{to command, order} castra parva facere. Alius mīles aquam attulit, alter parvum \voc[dōnum]{dōnum}{gift} parāvit: flōrem pulchrum ante avem posuit. “Hoc dōnum avis nōn intellegit,” puer in scholā dīxit, “sed causa bona est.” Magister eum laudāvit et \voc[causa]{causam}{cause, reason} in tabulā scrīpsit.

Interim hostēs per agrum currēbant. \voc[imperātor]{Imperātor}{commander, general, emperor} Rōmānus mīlitēs monuit: “Hostēs fugientēs nōn semper timendī sunt; interdum revertuntur.” Deinde mīlitēs Rōmānī hostēs \voc[fugō]{fugāvērunt}{to put to flight}. Dux tamen nūllum hostem \voc[occīdō]{occīdere}{to kill} cupīvit, quia bellum nōn amābat et vēram pācem quaerēbat.

Sub vesperum mīles avem iterum quaesīvit et tandem \voc[reperiō]{repperit}{to find}. Avis iam melius valēbat. Mīles gaudēns dīxit: “Avis vīvit!” Dux respondit: “Bene. Nunc castra relinquere possumus.” Ita mīlitēs castra \voc[relinquō]{relīquērunt}{to relinquish, leave} et ad urbem rediērunt.

In urbe erat \voc[theātrum]{theātrum}{theatre}. Ibi poēta fābulam dē ave vīvā et mīlite currente nārrāvit. \voc[vōx]{Vōx}{voice} poētae magna erat, et populus audiēns tacuit. Poēta dīxit: “Mora interdum bona est, sī causa bona est.” Tum magister discipulīs explicāvit: “\voc[mora]{Mora}{delay} dūcis avem servāvit.”

Post fābulam puella \voc[flōs]{flōrem}{flower} in librō pictum vīdit et dīxit: “Flōs est parvus, sed dōnum pulchrum est.” Puer addidit: “Et fābula est \voc[vērus]{vēra}{true, real}?” Magister respondit: “Fortasse nōn omnīnō vēra est; sed sententia vēra est.”

Tum magister novum exemplum scrīpsit: “Hoste fugātō, mīlitēs rediērunt. Ave repertā, dux laetus fuit. Locō relictō, lēgātus ad imperātōrem vēnit.” Discipulī intellexērunt participia nōn solum verba ornāre, sed etiam tempora et causās ostendere.

Ad finem lēctiōnis magister dīxit: “\voc[novus]{Nova}{new} verba hodiē habēmus, sed vetera verba quoque repetimus. \voc[cēterus]{Cēterī}{the other, rest} discipulī crās exempla scrībent.” Puer rīdēns respondit: “Magister, participia difficilia sunt, sed fābula bona est.” Magister ait: “Sī bene legitis, difficultās fugātur.”

\end{readersection}

\begin{readersection}{Rēgīna et Nāvis}
\voc[prīmum]{Prīmum}{first, at first} magister discipulīs fābulam dē urbe antīquā nārrāvit. In fābulā erat magnum \voc[forum]{forum}{forum, public place}; in forō stābat \voc[rēgīna]{rēgīna}{queen}. Rēgīna nōn erat superba, sed sapiēns; cīvēs eam amābant, quia iūstitiam et pācem servābat.

Prope forum erat magna \voc[porta]{porta}{gate}. Hostēs portam dēlēre cupīvērunt, sed porta ā cīvibus servāta est. Tum rēgīna ad populum \voc[loquor]{locūta est}{to say, speak}: “Cīvēs, nōlīte timēre. Urbs nostra nōn dēlēta est; porta servāta est; spēs nōbīs manet.”

In portū stābat \voc[nāvis]{nāvis}{ship}. Nauta pauper in nāvī sedēbat et pecūniam nōn habēbat. Is erat \voc[pauper]{pauper}{poor}, sed nōn \voc[stultus]{stultus}{foolish, stupid}. Rēgīna nautam vīdit et dīxit: “Hic vir miser est, sed bonus esse potest.” Tum nauta rēgīnae respondit: “Ego \voc[miser]{miser}{poor, miserable} sum, sed urbem adiuvāre cupiō.”

Nauta rēgīnae auxilium \voc[polliceor]{pollicitus est}{to promise}. “Nāve meā ūtar,” inquit, “et frumentum ad urbem feram.” Rēgīna nautam laudāvit et parvam \voc[pecūnia]{pecūniam}{money, wealth} eī dedit. Nauta autem dīxit: “Pecūnia mihi grāta est, sed \voc[glōria]{glōria}{glory, honor} cīvitātis mihi cārior est.”

\voc[dum]{Dum}{while} nauta nāvem parat, hostēs iterum ad portam veniunt. Ūnus mīles Rōmānus hostēs sequi \voc[audeō]{ausus est}{to dare, venture}. Aliī mīlitēs eum \voc[sequor]{secūtī sunt}{to follow}. Rēgīna eōs mīrāta est, quia paucī virī contrā multōs hostēs stāre ausī erant.

Hostēs forum occupāre \voc[cōnor]{cōnātī sunt}{to try, attempt}, sed ā cīvibus prohibiti sunt. Porta ā mīlitibus defensa est, et hostēs fugātī sunt. Tum rēgīna dīxit: “Bellum nōndum fīnītum est, sed animus cīvium magnus est.” Puer in scholā rogāvit: “Magister, quid est prohibiti sunt?” Magister respondit: “Hoc est perfectum passīvum.”

Interim nauta nāve \voc[ūtor]{ūsus est}{to use, employ} et frumentum in urbem tulit. Populus in forō eum exspectābat. Rēgīna dīxit: “Hic nauta nōn sōlum pecūniā, sed animō propriō urbem servāvit.” Magister discipulīs mōnstrāvit: “\voc[proprius]{Proprius}{one’s own, personal} cum nōmine convenit: animus proprius, pecūnia propria, dōnum proprium.”

Post proelium hostēs multa aedificia \voc[dēleō]{dēlēverant}{to destroy}, sed forum nōn dēlētum erat. Porta servāta erat, nāvis inventa erat, frumentum allātum erat. Rēgīna tum bellum \voc[fīniō]{fīnīvīt}{to finish}; pāx facta est, et cīvēs in forō gaudēbant.

Sed nōn omnēs laetī erant. Ūnus mīles in proeliō graviter vulnerātus erat et mox \voc[morior]{mortuus est}{to die}. Rēgīna eum laudāvit: “Ille vir pro urbe mortuus est; glōria eius nōn morietur.” Cīvēs tacitī stābant, quia mors mīlitis eōs movit.

Tum magister aliam partem fābulae legit: “Post multōs annōs puer in eādem urbe \voc[nāscor]{nātus est}{to be born}. Puer nōn rēx, nōn dux, sed cīvis bonus fuit. Is saepe dīxit: ‘Rēs \voc[cīvīlis]{cīvīlis}{civic, civil} nōn est parva; cīvitās omnibus cīvibus commūnis est.’”

Discipulī sententiam scrīpsērunt: “Urbs servāta est; porta defensa est; bellum fīnītum est; mīles mortuus est; puer nātus est.” Magister dīxit: “Ita perfectum passīvum intellegitis. Sed vidēte etiam verba quae fōrmam passīvam habent, sed sententiam actīvam: locūtus est, cōnātus est, secūtus est, ūsus est.”

Ad finem lēctiōnis puer rogāvit: “Magister, cūr nauta pecūniam nōn prīmum quaesīvit?” Magister respondit: “Quia nauta stultus nōn erat. Pecūnia utilis est, sed glōria honesta et salūs urbis interdum pluris valent.” Puella addidit: “Et rēgīna miserō nautā crēdidit.” Magister subrīsit: “Bene dīcis; fides cīvilis urbem servat.”

\end{readersection}

\begin{readersection}{Amīcitia Aurō Cārior}
Magister hodiē discipulīs dīxit: “In hāc lēctiōne comparātīva discimus: cārior, firmior, facilius, altius.” Deinde fābulam dē puerō et amīcō nārrāvit. Puer \voc[aliquandō]{aliquandō}{once, at one time} ad \voc[flūmen]{flūmen}{river} vēnit; aqua flūminis erat lata, sed via ad pontem nōn erat difficilis.

Cum puer ad flūmen stāret, amīcus eius \voc[citō]{citō}{quickly, fast} cucurrit. In manū habēbat parvum \voc[gladius]{gladium}{sword}; gladius erat ex \voc[ferrum]{ferrō}{iron}, nōn ex \voc[aurum]{aurō}{gold}. Puer rīsit et dīxit: “Ferrum utilius est quam aurum, sī perīculum adest.” Amīcus respondit: “Sed pāx melior est quam gladius.”

Sub arbore erat senex \voc[humilis]{humilis}{humble}. Is puerōs audīvit et dīxit: “\voc[amīcitia]{Amīcitia}{friendship} est melior quam \voc[dīvitiae]{dīvitiae}{riches, wealth}. Aurum saepe fāmam parit, sed amīcitia vītam firmiōrem facit.” Puer rogāvit: “Estne fāma melior quam amīcitia?” Senex respondit: “Nōn. \voc[fāma]{Fāma}{fame, reputation} sine virtūte levis est.”

Tum puer alterum amīcum vīdit. Is pallidus erat, quia \voc[morbus]{morbus}{disease, illness} eum vexābat. Medicus, vir doctus, ad eum vēnit et dīxit: “Morbus acer est, sed animus tuus firmus est.” Puer aquam dedit, et amīcus melius valuit. Magister discipulīs dīxit: “Vidē\-tisne? \voc[ācer]{Ācer}{sharp, vigorous, violent} morbus peior est quam labor, sed amīcitia morbum leviōrem facit.”

Postea omnēs ad collem \voc[altus]{altum}{high, tall} ascendērunt. Collis altior erat quam mūrus oppidī, et via \voc[lātus]{lātior}{wide, broad} erat quam via urbis. Super colle erat \voc[lūx]{lūx}{light} alba sōlis; puer dīxit: “\voc[albus]{Alba}{white} lūx pulchrior est quam aurum.” Senex respondit: “Et vēritās clārior est quam lūx.”

\voc[tum]{Tum}{then} magister in fābulā errōrem mōnstrāvit. Puer dīxerat aurum esse optimum, sed senex eum docuit. “Hic erat \voc[error]{error}{error, mistake} puerī,” inquit magister. “Aurum nōn est summum bonum; amīcitia est bonum firmius. Fāma interdum minor est quam virtūs, dīvitiae saepe minōrēs sunt quam sapientia.”

Ad flūmen iter facere nunc \voc[facilis]{facilius}{easy} erat, quia puerī cōnsilium melius habēbant. Amīcus gladium in vāgīnā tenuit et dīxit: “Ferrum ācrius est quam lignum, sed animus firmior est quam ferrum.” Senex laudāvit eum: “Bene dīcis. \voc[firmus]{Firmus}{firm, strong} animus est melior quam gladius ācer.”

In ripā flūminis puer parvam rem invēnit: signum vetus ex aurō. Omnēs signum spectāvērunt, sed senex dīxit: “Nōlīte dīvitias nimis amāre. Hoc aurum est minus ūtile quam amīcitia.” Puer signum magistrō dedit; magister respondit: “Hoc dōnum honestius est quam fāma.”

Post lēctiōnem discipulī sententiam scrīpsērunt: “Amīcitia aurō cārior est; virtūs fāmā maior est; animus firmus ferrō fortior est.” Magister addidit: “Hae sententiae sunt \voc[nātūrālis]{nātūrālēs}{natural}; nam homō nātūrā amīcitiam quaerit.” Discipulī verba comparātīva notāvērunt.

\voc[posterus]{Posterō}{following, next} diē puer ad magistrum rediit et dīxit: “Heri intellexī amīcitiam esse meliōrem quam aurum. Etiam error meus mihi prōfuit.” Magister rīsit: “Error interdum magister est. Sed discipulus sapiēns errorem \voc[minor]{minōrem}{smaller, less} facit, quia bene audit et discit.” Puer respondit: “Et amīcus bonus errōrem faciliōrem facit.”

\end{readersection}

\begin{readersection}{Hinc Hūc Itur}
Magister hodiē in scholā dīxit: “Nunc verbum īrregulāre discimus: \voc[eō]{eō}{to go}. Ego eō, tū īs, is it; nōs īmus, vōs ītis, eī eunt.” Discipulī verba audiunt et scrībunt. Puer ridet et dīcit: “Magister, hoc verbum breve est, sed nōn facile est.”

Tum magister fābulam narrat. In fābulā est \voc[senex]{senex}{old person}; cum sene est \voc[soror]{soror}{sister} eius. Senex vultum tranquillum habet, sed soror eius celeriter ambulāre cupit. Magister dīcit: “Vidēte \voc[vultus]{vultum}{look, countenance} senis: tranquillus est. Sed soror ad urbem īre cupit.”

Soror senī dīcit: “Frāter, \voc[hinc]{hinc}{hence, from this place} īmus; \voc[hūc]{hūc}{hither, to this place} crās redīmus.” Senex respondet: “Nōlī tam celeriter īre. \voc[aetās]{Aetās}{lifetime, age} mea longa est, et iam lentē eō.” Soror rīdet: “Sed iter breve est, et urbs prope est.”

Duo puerī viam mōnstrant. Ūnus dīcit: “Si ad portam īre vultis, hāc viā īte.” Alter autem monet: “Nōlīte per agrum īre, quia ibi mīlitēs sunt et \voc[arma]{arma}{arms, weapons} portant.” Senex arma videt et dīcit: “Arma saepe \voc[iniūria]{iniūriam}{injury, injustice} pariunt; pāx melior est.”

Soror tamen ad urbem \voc[adeō]{adit}{to go to, approach}. Senex lentē sequitur. Ad portam veniunt, et soror rogat: “Quō nunc īmus?” Senex respondet: “Prīmum forum adīmus; \voc[deinde]{deinde}{afterwards, then} domum redīmus.” In forō populus stat, et magistrātus cīvēs audit.

In forō rēx antīquus urbem \voc[condō]{condidisse}{to build, establish} dīcitur. Magister discipulīs explicat: “Rōmulus Rōmam condidit. Sed in hāc fābulā senex nōn urbem condit; tantum fābulam dē urbe audit.” Tum puer scrībit: “Rēx urbem condidit; senex urbem adit.”

Soror librum parvum ē sacculō capit et dīcit: “Hunc librum emere \voc[cupiō]{cupiō}{to desire}.” Senex respondet: “Ego librum legere cupiō, sed pecūniam servāre quoque cupiō.” Librārius librum mōnstrat, et soror eum accipit. Senex vultum laetum habet, quia soror tandem tacet.

Tum senex \voc[cōgitō]{cōgitat}{to think, intend}: “Tempus fugit; unus \voc[mēnsis]{mēnsis}{month} abit, alter mēnsis venit. Sed librī manent.” Soror audit et dīcit: “Nōlī nimis cōgitāre; nunc domum īmus.” Senex respondet: “Bene. Sed ante portam iterum īre volō.”

Ad portam redeunt. Soror ē forō \voc[exeō]{exit}{to go out, exit}; senex quoque exit. Puerī eōs vident et rogant: “Quō ītis?” Soror respondet: “Domum redīmus.” Senex addit: “Hinc exīmus, hūc crās redībimus.” Puerī rīdent, quia senex verba magistrī bene tenet.

In viā soror novum cōnsilium \voc[ineō]{init}{to go into, enter, begin}: “Crās iter longius faciemus.” Senex vultum gravem facit et dīcit: “Nōn cupiō. Ego \voc[tantum]{tantum}{only} ad forum īre possum.” Soror tamen \voc[optō]{optat}{to wish for, desire} ut frāter fortior sit. Senex respondet: “Optāre facile est; īre difficilius est.”

Tandem domum \voc[redeō]{redeunt}{to go back, return}. Senex sedet, soror librum legit, et magister fābulam finit. Discipulī intellegunt: hinc īmus, hūc redīmus; ex urbe exīmus, ad urbem adīmus; cōnsilium inīmus, domum redīmus. Verbum eō parvum est, sed multa itinera aperit.

\end{readersection}

\begin{readersection}{Philosophus et Scrīptor}
\voc[cōnstat]{Cōnstat}{it is well known} multōs discipulōs librōs legere velle. Hodiē magister fābulam dē \voc[philosophus]{philosophō}{philosopher} et \voc[scrīptor]{scrīptōre}{writer, author} nārrat. Philosophus in forō sedet et librum tenet; scrīptor ad eum venit, quia verba sapientiae audīre \voc[volō]{vult}{to want, wish}.

Scrīptor philosophō librum \voc[offerō]{offert}{to present, offer}. “Hunc librum tibi dare volō,” inquit. Philosophus librum capit, sed respondet: “Ego librum accipere volō, sed falsam glōriam nōlō.” Scrīptor rīdet: “Nōn glōriam quaerō; vēritātem trādere volō.”

Prope eōs stat \voc[mulier]{mulier}{woman}. Ea librum auditum vult, sed disputāre \voc[nōlō]{nōlit}{to not want}. Philosophus eam rogat: “Cūr tacēre māvīs?” Mulier respondet: “Quia verba sapientium saepe difficilia sunt, et ego errāre \voc[metuō]{metuō}{to fear}.” Philosophus dīcit: “Nōlī metuere; qui discere vult, rogāre dēbet.”

Tum scrīptor sententiam in tabulā scrībit: “\voc[fīnis]{Fīnis}{end, limit} discendī nōn est glōria, sed vēritās.” Philosophus laudat sententiam et dīcit: “Bene scrībis. Ego tamen \voc[putō]{putō}{to ponder, consider} fīnem sapientiae nōn esse facilem.” Mulier addit: “Ego scīre volō, sed multa adhūc \voc[nesciō]{nesciō}{to not know}.”

Philosophus tum varias sententias discipulīs mōnstrat. “Alius scrīptor glōriam vult; alius pecūniam māvult; alius pācem et sapientiam quaerit. Sed bonus scrīptor vēritātem \voc[trādō]{trādere}{to hand down, narrate} vult.” Discipulī audiunt et verba in librō scrībunt.

Mulier librum spectat et dīcit: “Haec verba nōn sunt \voc[aliēnus]{aliēna}{foreign, alien}; mihi placent. Nunc legere volō, nōn tacēre.” Scrīptor gaudet, quia mulier iam nōn metuit. Philosophus autem monet: “Nōlīte omnia statim crēdere. Legere vultis, sed etiam putāre dēbētis.”

Tum puer philosophum rogat: “Magister, utrum mālō librum legere an fābulam audīre?” Philosophus respondet: “Sī bene discere cupis, utrumque bonum est. Sed ego ipse legere \voc[mālō]{mālō}{to prefer}, quia librum iterum atque iterum inspicere possum.” Puer respondet: “Ego audīre mālō, sed crās legere volō.”

Ad fīnem lēctiōnis scrīptor dīcit: “Ego \voc[spērō]{spērō}{to hope} discipulōs hanc fābulam amātūrōs esse.” Philosophus rīdet: “Et ego spērō eōs nōn sōlum verba, sed etiam sententiam intellectūrōs esse.” In tabulā sunt \voc[varius]{varia}{various, different} verba: volō, nōlō, mālō. Discipulī intellegunt: aliud velle possumus, aliud nōlle, aliud mālle.

\end{readersection}

\begin{readersection}{Adulēscēns quī Stellam Vidit}
\voc[aestās]{Aestās}{summer} erat, quae urbem calidam faciēbat. Magister discipulīs fābulam nārrāvit dē puerō \voc[adulēscēns]{adulēscentī}{young}, quī cum patre ad flūmen proficīscī cupīvit. Pater eī dīxit: “Iter, quod hodiē facimus, breve est, sed verba, quae discimus, longa manent.”

Adulēscēns, quī librum in manū tenēbat, ad agrum vēnit. Ibi sub arbore senex \voc[sedeō]{sedēbat}{to sit, stay}, cui puer librum mōnstrāvit. Senex dīxit: “Liber, quem tenēs, bonus est. Sed nōlī sine cōnsiliō ambulāre, quia viā saepe \voc[errō]{errant}{to wander, err} eī quī nōn attendunt.”

Tum \voc[accidit]{accidit}{it happens} ut avis ē silvā subitō exīret. Avis parva, quam puer vīdit, mīlitem \voc[terreō]{terruit}{to frighten, terrify}, quia mīles sonum audīvit et hostem putāvit. Sed pater rīsit: “Haec avis, quae tē terret, nōn est hostis. Timor saepe rem parvam magnam facit.”

Prope viam stābat vir \voc[inimīcus]{inimīcus}{hostile, inimical}, quī puerō malum cōnsilium dedit. Is dīxit: “Nōlī patri crēdere; ego viam meliōrem sciō.” Sed senex, cuius sapientia omnibus nota erat, respondit: “Cōnsilium, quod dat, \voc[falsus]{falsum}{false, untrue} est. Vir, quī falsum cōnsilium dat, amīcus nōn est.”

Adulēscēns tamen nōn timuīt. Is lapidem parvum in aquam \voc[iaciō]{iēcit}{to throw, hurl, cast}. Aqua, quam lapis tetigit, parvās undās \voc[gignō]{genuit}{to produce, yield}. Magister hīc discipulīs dīxit: “Vidētisne? Causa parva saepe effectum maiōrem gignit. Sententia, quam discimus, utilis est.”

Pater fīliō parvum dōnum dedit: tabulam, in quā \voc[stella]{stella}{star} picta erat. “Haec stella,” inquit, “tē admoneat nocte, quā viam quaeris.” Adulēscēns dōnum accēpit et patri grātiās ēgit. Dōnum, quod pater dedit, puerō \voc[iūcundus]{iūcundum}{pleasant, agreeable} erat.

Posteā puer et pater ad urbem \voc[proficīscor]{profectī sunt}{to depart, leave}. In viā puer fessus erat, sed iter \voc[levis]{leve}{light, swift} vidēbātur, quia pater eum verbīs bonīs \voc[iuvō]{iuvābat}{to help, please}. Pater dīxit: “Amor eōs iuvat quī labōrant; verba bona dolorem minuunt.”

Sub vesperum puer parvum \voc[dolor]{dolōrem}{pain, sorrow} in pede sensit. Pater eum adiūvit et dīxit: “Dolor, quem sentīs, levis est; crās bene valēbis.” Puer respondit: “Pater, verba tua mē iuvant.” Tum senex, quī eōs sequēbātur, addidit: “Dolor, quī facile fertur, animum docet.”

Cum ad urbem rediērunt, pater librum puerō \voc[reddō]{reddidit}{to give back, return}. “Hic liber,” inquit, “quem mihi herī dedistī, nunc tibi redditur.” Puer librum accēpit et in forō legit. Populus, quī verba audīvit, puerum laudāvit; sed puer dīxit: “\voc[laus]{Laus}{praise} patri et magistrō debētur, quī mē docuērunt.”

Magister fābulam fīniēns discipulīs dīxit: “Sī sententiam intellegere vultis, attendite ad eum quī agit, ad rem quam videt, ad virum cui liber datur, ad patrem cuius verba prosunt. Ita relativa nōn difficilia sunt.” Discipulī responderunt: “Hoc \voc[facile]{facile}{easily} intellegimus, \voc[ut]{ut}{as, just as} tū nārrās.”

\end{readersection}

\begin{readersection}{Quid ad Rēm Pūblicam Pertineat}
Magister hodiē discipulīs dīxit: “Nunc discimus quōmodo coniūnctīvus in sententiīs positus sit. Interrogō ubi sit urbs; hortor ut audiātis; rēfert quid sciātis.” Discipulī verba in tabulā scrīpsērunt.

In fābulā erat \voc[senātus]{senātus}{senate} Rōmānus. Senātus dē urbe, dē bellō, dē pāce dēliberābat. Magister discipulōs rogāvit quid senātus ageret. Puer \voc[respondeō]{respondit}{to respond, reply}: “Senātus quaerit quid ad cīvitātem \voc[pertineō]{pertineat}{to concern, pertain to}.”

Tum nūntius vēnit et dīxit \voc[classis]{classem}{fleet} prope mare stāre. Cōnsul nesciēbat utrum classis amīca esset an inimīca. Itaque senātus quaesīvit ubi classis esset, quō īret, et cūr ad urbem venīret. Magister dīxit: “Ecce, discipulī: haec sunt interrogātiōnēs oblīquae.”

In eādem urbe vīvēbat vir \voc[dīves]{dīves}{rich, wealthy}, quī pecūniam multam habēbat. Sed vir dīves saepe \voc[negō]{negābat}{to deny, say not} sē cīvitātī pecūniam dare posse. Cīvēs tamen sciēbant eum pecūniam habēre. Magister interrogāvit utrum vir vērum dīceret an falsum negāret.

Illī virō erat \voc[fīlia]{fīlia}{daughter} sapiēns. Fīlia patrem \voc[hortor]{hortāta est}{to encourage, urge} ut senātuī responderet et ut pecūniam cīvitātī daret. Ea dīxit: “Pater, \voc[rēfert]{rēfert}{it matters, it is important} quid agās. Pecūnia tua nōn sōlum ad tē, sed etiam ad cīvēs pertinet.”

Pater autem respondit: “Ego \voc[īgnōrō]{īgnōrō}{to not know} quid faciendum sit.” Fīlia dīxit: “Nōlī id dīcere. Nōn difficile est scīre quid rēctum sit.” Tum magister discipulīs explicāvit: “Quid faciendum sit, quid rēctum sit: hīc quoque coniūnctīvus est.”

Senātus virum dīvitēm vocāvit. Cōnsul eum hortātus est ut aurum in mensā \voc[pōnō]{pōneret}{to put, place}. Vir prīmum tacuit; deinde aurum posuit. Fīlia laeta fuit, quia pater tandem intellexit quid ad rem pūblicam pertineret.

Posteā senātus statuit ut \voc[lūdus]{lūdī}{game, sport} fierent, nōn glōriae causā, sed pācis causā. Puerī, puellae, mīlitēs et nautae in forō erant. Magister dīxit: “\voc[causā]{Causā}{for the sake of, on account of} cum genitīvō ponitur: pācis causā, glōriae causā, lūdōrum causā.”

In lūdīs puer \voc[sānus]{sānus}{sound, healthy} currēbat, sed senex currere nōn poterat. Fīlia virī dīvitis senem iuvit. Cōnsul rogāvit cūr eum adiuvāret. Ea respondit: “Quia corpus senis nōn est firmum, sed animus eius sānis cīvibus similis est.” Senātus hanc sententiam laudāvit.

Tum magister discipulīs dīxit: “Sententia \voc[rēctus]{rēcta}{straight, right} saepe facilis nōn est. Sed rēfert ut eam quaerāmus.” Puer rogāvit: “Magister, estne hoc \voc[facilis]{facile}{easy}?” Magister respondit: “Nōn semper facile est; sed qui bene legit, facilius intellegit.”

Sub finem fābulae cōnsul spectāvit \voc[orbis]{orbem}{circle, world} terrae in tabulā pictum. “Orbis magnus est,” inquit, “sed urbs nostra parva est. Tamen etiam parva urbs iūstitiam habēre potest.” Senātus respondit oportēre ut cīvēs rēctē agerent et ut pācem servārent.

Magister lēctiōnem fīnīvit: “Hodiē didicimus quid senātus quaereret, quid cōnsul iuberet, quid fīlia hortārētur, quid vir dīves negāret. Nunc intellegitis coniūnctīvum nōn sōlum formam esse, sed etiam modum cogitandī.” Discipulī responderunt sē rem intellegere.

\end{readersection}

\begin{readersection}{Nē Amor Pereat}
Magister hodiē discipulīs dīxit: “In hāc lēctiōne coniūnctīvus in sententiā prīncipālī discitur. Dīcere possumus: vīvat patria; nē urbs pereat; taceāmus; quid faciam?” Tum fābulam dē urbe parvā nārrāvit.

In urbe erat \voc[nox]{nox}{night}. Cīvēs in forō stābant, quia hostēs prope mūrōs erant. Cōnsul dīxit: “Nē urbs \voc[pereō]{pereat}{to perish, be ruined}! Cīvēs, portas servēmus et auxilium petāmus.” Populus tacitus erat, sed timor multōs tenēbat.

Tum senex, quī multōs annōs in urbe vīxerat, prōcessit et dīxit: “Nē timōrī \voc[cēdō]{cēdāmus}{to go, yield to}. Sī timōrī cēdāmus, cīvitās perībit. Sed sī fidēs et \voc[amor]{amor}{love} maneant, spēs quoque manēbit.” Hic \voc[sermō]{sermō}{speech, talk, discourse} cīvēs movit.

Puella parva patrem rogāvit: “Quid faciamus?” Pater respondit: “\voc[auxilium]{Auxilium}{help, aid} parēmus. Aquam ferāmus, portās claudāmus, puerōs et senēs servēmus.” Tum puella dīxit: “Ego quoque aquam feram.” Pater eam laudāvit: “Tū parva es, sed magnā grātiā digna.”

Cōnsul mīlitēs monuit: “Nē hostēs urbem intrent. Hostēs ā portā \voc[prohibeō]{prohibeantur}{to prevent, prohibit}. Nē cīvēs pāce \voc[careō]{careant}{to lack}. Nē populus spem amittat.” Deinde ad cīvēs conversus est: “Quisque suam \voc[pars]{partem}{part} agat. Alius portam servet, alius aquam ferat, alius vulnerātōs iuvet.” Ita cīvēs officia inter sē divīsērunt.

Interim mīles iuvenis dīxit: “Ego hostēs impetū petam.” Senex respondit: “Nē īrā vincāris. Fortis esto, sed prūdentiam tene. Nōn glōria, sed salūs urbis quaerenda est.” Cōnsul hoc cōnsilium bonum \voc[exīstimō]{exīstimāvit}{to consider, estimate}.

Puella, quae aquam portābat, vulnerātum mīlitem vīdit. “Utinam medicus hīc adsit!” inquit. Tum medicus vēnit et mīlitī opem tulit. Mīles dīxit: “Tibi \voc[grātia]{grātiam}{grace, thankfulness} habeō.” Medicus respondit: “Nē mihi grātiās agās; urbs ipsa servanda est.”

Post breve proelium hostēs discessērunt. Cīvēs gaudēbant, sed cōnsul dīxit: “Nē nimis celeriter gaudeāmus. Portās \voc[iterum]{iterum}{again} inspiciāmus; mūrōs iterum videāmus; vigilēs \voc[parō]{parēmus}{to arrange, prepare, provide}.” Senex addidit: “\voc[sānē]{Sānē}{truly, indeed, really} prūdenter loqueris.”

Nox tamen longa erat, et puerī dormīre cupiēbant. Magister in fābulā dīxit: “Nunc cīvēs \voc[taceō]{taceant}{to be silent}. Sī taceant, hostēs nihil audient.” Omnēs igitur tacuērunt \voc[ūsque]{ūsque}{all the way} ad prīmam lūcem. Silentium cīvitātī prōfuit.

Mane cōnsul cīvēs in forō laudāvit: “Urbs servāta est, quia amor patriae vōs \voc[dēlectō]{dēlectat}{to delight, please}, nōn glōria vana. Omnēs laude \voc[dignus]{dignī}{worthy of} estis.” Tum puella subrīsit et dīxit: “Utinam semper pāx sit; utinam amor numquam pereat.” Magister lēctiōnem fīnīvit: “Hoc discāmus: nē timōrī cēdāmus, sed ratiōnem sequāmur.”

\end{readersection}

\begin{readersection}{Ut Puerī Nōn Dubitent}
Magister hodiē discipulīs dīxit: “In hāc lēctiōne discimus sententiās fīnālēs: veniō ut discam; moneō nē errēs. Ita cōnsilium et causa faciendī mōnstrantur.” Tum fābulam dē scholā, hortō et hieme nārrāvit.

\voc[hiems]{Hiems}{winter} erat, sed discipulī ad scholam vēnērunt ut discerent. Magister fenestram clausit nē frīgus puerōs impedīret. In scholā erat \voc[hōra]{hōra}{hour, time} lēctiōnis, et omnēs tacēbant ut verba magistrī audīrent.

Magister fābulam dē agricola legit. Agricola in agrō labōrābat ut \voc[frūctus]{frūctūs}{product, fruit} bonōs parāret. Sed hiems agricolam saepe \voc[impediō]{impediēbat}{to impede, hinder}; nix agrō \voc[obstō]{obstābat}{to thwart, hinder}, et agricola tamen nōn dubitābat. Magister puerōs monuit ut labor hiemis eōs nōn \voc[dēterreō]{dēterrēret}{to deter, prevent}. Is labōrābat ut familiam aleret.

Puer, quī prope magistrum sedēbat, \voc[dubitō]{dubitāvit}{to doubt, hesitate}: “Cūr agricola in hieme labōrat?” Magister respondit: “Labōrat ut postea frūctūs habeat. Nōlī dubitāre; saepe homō hodiē labōrat ut crās melius vīvat.” Tum in tabulā scrīpsit: labōrat ut habeat.

Subitō puer alter per fenestram avem audīvit. Avis \voc[canō]{canēbat}{to sing}, quamquam hiems erat. Puer surrexit ut avem vidēret, sed magister eum monuit nē lēctiōnem relinqueret. “Aperī \voc[auris]{aurēs}{ear},” inquit, “ut fābulam audiās; oculōs postea ad avem vertēs.”

Tum magister discipulīs \voc[modus]{modum}{manner, way} bene legendī mōnstrāvit. “Prīmum audīte,” inquit, “ut intellegātis; deinde legite ut verba teneātis; tum scrībite nē memoria pereat.” Discipulī hunc modum probāvērunt, quia facilior erat quam antea.

Post lēctiōnem magister parvum pānem et māla in mensā posuit. Puerī \voc[edō]{edērunt}{to eat}; sed magister dīxit: “Nōlīte nimis edere, nē somnus veniat.” Puer rīdēns respondit: “Edimus ut valeāmus, nōn ut dormiāmus.” Magister eum laudāvit.

Puella tamen \voc[dubium]{dubium}{doubt} habēbat. Ea rogāvit: “Magister, utrum veniō ut discam, an veniō nē nesciam?” Magister respondit: “Utrumque rēctē dīcitur. Veniō ut discam; veniō nē nesciam. Hoc est discrīmen inter ut et nē.” Puella iam melius intellexit.

Deinde magister brevem \voc[ōrātiō]{ōrātiōnem}{oration, speech} habuit dē \voc[hūmānitās]{hūmānitāte}{humanity}. “Discimus,” inquit, “nōn sōlum ut verba sciāmus, sed etiam ut hūmānitātem colāmus. Qui alterum audit, hūmānior fit. Qui alterum semper impedit, hūmānitātem minuit.”

In ōrātiōne magister saepe \voc[ōs]{ōre}{mouth} et aurēs memorāvit: “Ōs habēmus ut loquāmur; aurēs habēmus ut audiāmus. Sed sapiēns nōn semper loquitur; saepe tacet ut alium audiat.” Discipulī tacitī erant, quia sententia eīs placēbat.

\voc[interdiū]{Interdiū}{in the daytime} discipulī in scholā legēbant; \voc[noctū]{noctū}{at night} domī verba repetēbant. Puer dīxit: “\voc[adhūc]{Adhūc}{so far, still} multa nesciō, sed nōn dubitō. Legō ut sciam, scrībō nē oblīvīscar, audiō ut intellegam.” Magister respondit: “Bene dīcis; hic modus discendī bonus est.”

Ad fīnem hōrae magister discipulōs monuit nē labōrem timērent. “Hiems transit,” inquit, “et frūctus post labōrem venit. Discite ut vōbīs prōsitis; scrībite ut verba maneant; canite interdum ut animus laetus sit.” Tum puerī parvum carmen canere coepērunt, et lēctiō fīnīta est.

\end{readersection}

\begin{readersection}{Cum Litterae Perditae Essent}
Magister hodiē discipulīs dīxit: “In hāc lēctiōne cum cum coniūnctīvō discimus. Cum puer legat, magister audit; cum litterae perditae sint, omnēs quaerunt; cum discipulus errāverit, tamen discere potest.”

In fābulā puer \voc[littera]{litterās}{letter} magistrī portābat. Litterae dē \voc[geōmetria]{geōmetriā}{geometry} erant, quia magister in librō novum \voc[genus]{genus}{kind, class, race} figūrae discipulīs mōnstrāre volēbat. Sed cum puer per forum īret, litterās perdidit.

Cum puer litterās nōn invenīret, valde timuit. Is dīxit: “Ego \voc[peccō]{peccāvī}{to make a mistake, sin}; litterās \voc[perdō]{perdidī}{to lose, destroy}.” Tum soror eius, puella bona, respondit: “Nōlī tacēre. Cum verum dīcās, magister fortasse tē nōn damnābit.”

Puer ad scholam vēnit et magistrō omnia nārrāvit. Magister prīmum tacuit; deinde dīxit: “\voc[licet]{Licet}{it is permitted} errāre, sed nōn licet mentīrī. Cum litterās perdideris, novās litterās scrībere dēbēs.” Puer respondit: “Magister, poenam \voc[patior]{patī}{to suffer, endure} possum.”

Sed magister vir \voc[iūstus]{iūstus}{just} erat. Is puerum nōn expulit, sed eī tabulam dedit. “Scrībe,” inquit, “quid acciderit. Cum culpam tuam intellegās, iam melius discis.” Puer tabulam \voc[accipiō]{accēpit}{to accept, receive} et statim scrībere coepit.

Interim discipulus alius dīxit: “Puer ob culpam suam \voc[exilium]{exilium}{exile} meret.” Magistra autem respondit: “Nōn ita. Exilium grave est. Parvum \voc[vitium]{vitium}{vice, crime} nōn semper magnā poenā pūniendum est.” Cum hoc dīxisset, discipulī tacuerunt.

Tum magister cōnsilium \voc[cōnstituō]{cōnstituit}{to set, establish, decide}: “Puer litterās perditās iterum scrībat; cēterī eum iuvent, nē labōre nimis premātur.” Puella dīxit: “Hoc cōnsilium est \voc[honestus]{honestum}{honourable, decent, honest}.” Magister respondit: “Et iūstum est.”

Post hōram puer novās litterās scrīpsit. Cum litterae in mēnsā essent, magister eas legit et laudāvit. “Nunc,” inquit, “geōmetriam discere possumus.” Discipulī figūrās spectāvērunt, et magister dīxit: “Hoc genus figūrae facile intellegitur, sī animus attentus est.”

Sed puer adhūc maestus erat \voc[ob]{ob}{on account of} errōrem suum. Magister eum monuit: “Nē tē ipsum \voc[expellō]{expellās}{to expel}. Qui errat sed verum dīcit, ex animō suō expellere dēbet timōrem, nōn spem.” Puer haec verba accēpit et melius valuit.

Sub finem lēctiōnis omnēs ex scholā \voc[ēgredior]{ēgressī sunt}{to go out}. Cum ēgrederentur, puer magistrō dīxit: “Grātiās tibi agō, quia mē nōn expulistī.” Magister respondit: “Ego iniūriam \voc[dēfendō]{dēfendere}{to defend} nōn volō; sed discipulum honestum servāre volō.”

Deinde magister addidit: “Tamen hoc memineris: \voc[ultimus]{ultimus}{last, final} errōr nōn erit, nisi semper cūram adhibeās. Cum neglegenter agās, litterās iterum perdes.” Puer rīsit et dīxit: “Nōn iterum peccābō.” Magister respondit: “Nōn omnia \voc[vetō]{vetō}{to forbid, prevent, oppose}; sed mentīrī semper vetō.”

Discipulī domum rediērunt, et puer litterās novās portāvit. Cum nox venīret, ille in librō scrīpsit: “Hodiē peccāvī, sed verum dīxī. Magister adeō bonus fuit, ut spem nōn perderem.” Ita puer intellexit errōrem posse magistrum esse, cum animus honestus maneat.

\end{readersection}

\begin{readersection}{Dum Herba Crēscit}
Magister hodiē discipulīs dīxit: “Dum legimus, discimus; dum exspectāmus, animus saepe movētur; dum bene agimus, vīta crēscit.” Tum fābulam dē agrō, equō et plēbe nārrāre incipit.

In agrō erat \voc[equus]{equus}{horse} pulcher. Prope equum sedēbat \voc[eques]{eques}{horseman, cavalryman}, quī puerōs spectābat dum eī per viam ambulant. Equus \voc[herba]{herbam}{herb} ēdēbat, et herba sub sole \voc[crēscō]{crēscēbat}{to grow, increase}. Puer dīxit: “Dum herba crēscit, equus gaudet.”

Sed eques nōn gaudēbat. Is nūntium \voc[exspectō]{exspectābat}{to await}, dum sōl altius ascenderet. Nūntius tandem vēnit et dīxit: “In urbe \voc[plēbs]{plēbs}{common people, plebeians} auxilium petit.” Eques respondit: “Si plēbs auxilium petit, nōs \voc[ops]{opem}{power, strength, ability} ferre dēbēmus.”

Tum eques iter \voc[incipiō]{incipit}{to begin, undertake}. Equum movet et ad urbem contendit. Dum equus currit, pulvis surgit; dum puerī spectant, magister dīcit: “Vidētisne? Dum hoc agitur, alia quoque fiunt. Hoc dum tempus simul mōnstrat.”

In viā iacet magnum \voc[saxum]{saxum}{stone}. Equus saxum timet et pēs eius paene laeditur. Eques descendit, saxum \voc[moveō]{movet}{to move, set in motion}, et iter rūrsus incipit. “Equum servāre \voc[decet]{decet}{it is proper},” inquit, “quia equus amīcus est hominis.”

Mox eques ad urbem pervenit. Ibi \voc[reus]{reus}{defendant, accused} ante magistrātum stat. Plēbs clāmat, sed magistrātus tacet dum reus respondet. Reus dīcit: “Nōn sum \voc[barbarus]{barbarus}{barbarian, foreign}; legēs intellegō et iūstitiam petō.” Magister discipulīs dīcit: “Dum reus loquitur, iūdex audīre dēbet.”

In forō quoque stat philosophus. Is dē \voc[philosophia]{philosophiā}{philosophy} cum iuvenibus loquitur: “Dum homō philosophiae studet, animus eius crēscit. Sed philosophia nōn verbīs sōlīs, sed vītā probātur.” Puer hoc audit et sententiam scrībit.

Magistrātus tandem cōnsilium \voc[īnstituō]{īnstituit}{to establish, set up, start on}: “Prīmum reum audiāmus; deinde testēs petāmus; postrēmō iūdicium faciamus.” Plēbs tacet, dum cōnsilium īnstituitur. Eques quoque tacet, quia iūstitia sine silentio difficilis est.

Interim puerī in forō \voc[lūdō]{lūdunt}{to play}. Magister eōs monet: “Nōlīte lūdere dum iūdex loquitur.” Puerī respondent: “Lūdimus dum exspectāmus.” Magister ait: “Exspectāre licet, sed iūdicium honōrāre decet.” Tum puerī lūdum relinquunt.

Reus magistrātum \voc[petō]{petit}{to seek, ask for}: “Iūstitiam petō, nōn grātiam.” Eques addit: “Dum vēritās quaeritur, nemo damnārī dēbet.” Haec verba plēbem movent. Magistrātus reum audit, testēs audit, et tandem dīcit: “Reus nōn damnātur.”

Post iūdicium plēbs gaudet, eques equum ad aquam dūcit, et puerī rūrsus lūdunt. Eques spectat dum equus aquam bibit. “Dum aqua bibitur,” inquit, “equus melius valet.” Puer respondet: “Et dum iūstitia servātur, cīvitās melius valet.”

Magister lēctiōnem fīnit: “Hodiē didicimus: dum herba crēscit, equus ēst; dum eques nūntium exspectat, plēbs auxilium petit; dum reus loquitur, iūdex audit; dum puerī exspectant, lūdunt. Sed etiam didicimus: dum iūstitia petitur, cīvitās crēscit.”

\end{readersection}

\begin{readersection}{Qui Laborī Serviat}
Magister hodiē discipulīs dīxit: “In hāc lēctiōne discimus relātīvās sententiās cum coniūnctīvō. Possumus dīcere: quaerō virum quī labōrem ferat; dignus est quī laudētur; locus est unde aqua petātur.”

Tum fābulam dē fabrō nārrāvit. In oppidō erat vir quī templum reficere vellet. Ante templum \voc[iaceō]{iacēbat}{to lie, be situated} \voc[māteria]{māteria}{material}: lignum, saxum, ferrum. Saxum novum veterī saxō \voc[similis]{simile}{similar} erat, ut templum idem viderētur. Vir dīxit: “Hīc satis māteriae est, ex quā novum templum fiat.”

Populus magistrātum rogāvit ut virōs mitteret quī templō servirent. Magistrātus respondit: “Quaerimus hominēs quī \voc[labor]{labōrem}{labour, toil} nōn timeant.” Tum agricola, nauta, faber et discipulus vēnērunt. Omnēs dīxērunt sē labōrāre velle, ut templum servārētur.

Magistrātus eīs \voc[imperō]{imperāvit}{to command, rule}: “Lignum ferte, saxum movēte, viam aperīte.” At senex monuit: “Nē labōrī tantum imperēs. Oportet etiam virōs quaerere quibus labor nōn \voc[noceō]{noceat}{to hurt, injure}.” Magister discipulīs dīxit: “Vidētisne? Quibus noceat: relātīva sententia cum coniūnctīvō.”

Inter virōs erat servus quī diū dominō \voc[serviō]{servīverat}{to be a slave, serve}. Is tamen nōn erat miser animō; spem habēbat et templum sacrum amābat. “Ego,” inquit, “templō serviam, nōn sōlum dominō.” Populus eum laudāvit, quia labor eius omnibus prōfuit.

Sub pariete erat puer quī saxum tollere \voc[nequeō]{nequībat}{to be unable, cannot}. Faber eī dīxit: “Nōlī dolēre. Sunt labōrēs quōs puer facere nequeat, sed sunt etiam verba quae puer scrībat et memoriae trādat.” Puer igitur tabulās parāvit.

Senex discipulīs narrāvit \voc[memoria]{memoriam}{memory, remembrance} templī antīquī. “Hoc templum,” inquit, “per multa \voc[saeculum]{saecula}{generation, era, century} stetit. Memoria eius sacra est.” Puer rogāvit: “Estne templum \voc[sacer]{sacrum}{holy, sacred, divine} propter deōs?” Senex respondit: “Etiam propter hominēs quī fidē templum servāvērint.”

Tum terra sub mūrō movērī coepit. \voc[mōtus]{Mōtus}{motion} parvus erat, sed omnēs timuērunt. Faber dīxit: “Nōn est mōtus tantus ut mūrus cadat, sed satis magnus est quī nōs moneat.” Virī saxum firmāvērunt, nē templum perīret.

Magister hīc explicāvit: “Talis erat mōtus quī nōs monēret; dignī sunt virī quī laudentur; quaeritur locus \voc[unde]{unde}{from where, hence} aqua ferātur. Hae relātīvae sententiae nōn simplicem rem narrant, sed finem, effectum, causam aut proprietātem mōnstrant.”

Deinde sacerdos venit et dīxit: “Hoc opus \voc[iniūstus]{iniūstum}{unjust} nōn sit. Nē pauper labōret dum dīves spectet; nē servus solus labōrem ferat dum dominus taceat.” At dominus respondit: “Vera dīcis. Ego quoque labōrābō.” \voc[quidem]{Quidem}{indeed, certainly} dominus lentē labōrābat, sed labōrābat.

Puer tabulārum \voc[sēnsus]{sēnsum}{sensation, perception} intellegere cupīvit. Rogāvit: “Magister, quid est sēnsus huius fābulae?” Magister respondit: “Sēnsus est hic: labor, quī omnibus prōsit, nōn est servitūs, sed officium; et templum, quod memoriae serviat, populō quoque servit.”

Post multōs diēs templum restitūtum est. Populus dīxit: “Hoc opus est dignum quod memoriā teneātur.” Puer autem in tabulā scrīpsit: “Māteria parit formam, labor parit templum, memoria parit amōrem.” Magister addidit: “\voc[pariō]{Parere}{to bring forth, produce} nōn sōlum mātrum est; etiam labor fructum parit.”

\voc[oportet]{Oportet}{it is necessary, it is proper} igitur nōs quaerere labōrem quī prosit, nōn labōrem quī noceat. Oportet quoque memoriās servāre quae animōs meliorēs faciant. Ergō discipulī intellexērunt cūr relātīvae sententiae cum coniūnctīvō saepe nōn rem certam, sed causam, finem, effectum aut proprietātem significent.

\end{readersection}

\begin{readersection}{Ars Disputandī}
Magister hodiē discipulīs dīxit: “In hāc lēctiōne gerundium discimus: ars legendī, causa scrībendī, ad discendum, scrībendō discimus.” Tum tabulam parāvit et \voc[initium]{initium}{beginning, start} lēctiōnis fēcit.

In tabulā erat parvum \voc[signum]{signum}{sign, mark}. Magister colōrem signī mōnstrāvit et dīxit: “Hoc signum rubrum est; sed \voc[color]{color}{color} mūtārī potest. Ad colōrem mūtandum \voc[īnstrūmentum]{īnstrūmentum}{instrument, tool} parvum habēmus.” Puer īnstrūmentum capit, sed magister eum monet nē nimis celeriter agat.

Puer rogat: “Magister, cūr signum mūtāmus?” Magister respondet: “Ut rem clārius videāmus. Hāc ratiōne rem \voc[dēmōnstrō]{dēmōnstrāmus}{to demonstrate, prove}. Dēmōnstrandī causā signa interdum mūtanda sunt.” Puella sententiam scrībit: “Signum ad dēmōnstrandum prōdest.”

Deinde discipulī inter sē \voc[disputō]{disputant}{to debate, argue}. Ūnus dīcit colōrem rubrum esse pulchriōrem; alter colōrem album esse clāriōrem putat. Magister rīdet et dīcit: “Disputāre licet, sed cum \voc[prūdentia]{prūdentia}{prudence} disputandum est. Ars \voc[perītus]{perītī}{skillful, experienced} disputandī nōn clāmōre, sed argūmentīs probātur.”

Tum magister \voc[argūmentum]{argūmentum}{argument} bonum quaerit. Puer dīcit: “Color ruber melius vidētur.” Puella respondet: “Sed color albus facilius legitur.” Magister utrumque argūmentum laudat: “Haec \voc[sententia]{sententia}{opinion, meaning, sentence} bona est, illa quoque bona est. In disputandō sententiae audiendae sunt.”

Subitō aliquid \voc[ēveniō]{ēvenit}{to happen, occur}: puer īnstrūmentum in aquam cadere sinit. Omnēs rīdent, sed puer rubet. Magister dīcit: “Nōn est magnum malum. Errāre potest quisque; sed \voc[stultitia]{stultitia}{stupidity, foolishness} est errōrem tegere.” Puer respondet: “Mē \voc[paenitet]{paenitet}{it causes regret}; īnstrūmentum cūrāre dēbuī.”

Magister puerum nōn accūsat. “Nōn tē damnō,” inquit, “quia verum dīcis. Sed hinc discere dēbēs.” Deinde discipulīs imperat ut manūs lavent et tabulam tergant. “Disciplinam,” inquit, “cotīdiē \voc[exerceō]{exercēmus}{to exercise, practise}. Exercendō melius discimus.”

Post parvam moram magister discipulōs rogat quid didicerint. Puella respondet: “Didicimus colōrem mūtārī posse, signum ad dēmōnstrandum prōdesse, et argūmenta in disputandō necessāria esse.” Puer addit: “Didicī quoque mē īnstrūmentum neglegere nōn dēbēre.”

Magister tum ait: “Nōlīte \voc[ignōrō]{ignōrāre}{to not know} vim gerundiī. Ars legendī, modus scrībendī, causa discendī, īnstrūmentum pingendī: haec omnia in Latīnā linguā saepe inveniuntur.” Discipulī verba scrībunt, nē ea oblivīscantur.

Ad fīnem lēctiōnis magister \voc[cūnctus]{cūnctōs}{total, complete} discipulōs laudat: “Cūnctī hodiē bene labōrāvistis. Disputandō, scrībendō, audiendō et errōrēs corrigendō didicistis.” Tum puer, quī īnstrūmentum perdiderat, dīcit: “Crās prūdentius agam.” Magister respondet: “Hoc est bonum initium discendī.”

\end{readersection}

\chapter{Lectio provectior}
\begin{readersection}{Mercator, Virgo et Iudicium in Foro Romano}
\voc[hodie]{Hodie}{today} \voc[Romae]{Romae}{at Rome} in \voc[forum]{foro}{forum} \voc[quidam]{quidam}{a certain} \voc[emptor]{emptor}{buyer} stat. Is \voc[argentum]{argentum}{silver} in sacculo \voc[porto]{portat}{to carry} et \voc[vinum]{vinum}{wine} \voc[Sabinum]{Sabinum}{Sabine wine} quaerit. Tabernarius amphoras multas habet atque clamat: “\voc[ecce]{Ecce}{look}, bonum vinum! Hoc vinum \voc[verus]{verum}{true} est, non fictum.” Emptor tabernarium \voc[rogo]{rogat}{to ask}: “Quantum pretium est? Num pretium \voc[modicus]{modicum}{moderate} est?” Tabernarius respondet: “Pretium non est magnum; argentum tuum \voc[sufficio]{sufficit}{to be enough}.” Emptor tamen \voc[nego]{negat}{to deny} se statim emere velle, nam prius vinum gustare cupit.

Iuxta tabernam sedet \voc[femina]{femina}{woman} pauper, et cum ea est \voc[virgo]{virgo}{maiden} parva. Virgo capillos suos floribus \voc[como]{comit}{to adorn} et oculos ad caelum tollit, ubi unum \voc[sidus]{sidus}{star} etiam inter nubes lucet. Femina, quae multos annos in foro vendit, semper \voc[memor]{memor}{mindful} veteris fortunae est. Ea dicit: “O filia, si animum bonum \voc[servo]{servas}{to preserve}, eris \voc[beatus]{beata}{happy}. \voc[semper]{Semper}{always} pauca nobis adsunt, sed pauca interdum satis sunt.” Virgo ridet; \voc[animus]{animus}{mind} eius levis est, nec \voc[timor]{timor}{fear} eam torquet.

Tum advenit iuvenis \voc[pulcher]{pulcher}{beautiful}, qui virginem diu amat et eam ducere uxorem cupit. Femina dicit: “Filia mea nondum \voc[nubo]{nubit}{to marry}; sed si iuvenis iustus est, posteri dies rem probabunt.” Iuvenis respondet: “Ego nihil iniquum peto; hoc solum ad me \voc[pertineo]{pertinet}{to pertain}. Virginem amo et familiam vestram iuvare volo.” Emptor verba audit et putat iuvenem verum dicere. Sed tabernarius \voc[prior]{prior}{former} surgit et clamat: “Nolite credere! Iste iuvenis olim \voc[furtum]{furtum}{theft} fecit. Ego eum \voc[damno]{damno}{to condemn}.”

Iuvenis irascitur, sed virgo \voc[fleo]{flet}{to weep}. Mater eius dicit: “Ne flete; causam audiamus. Si crimen est, \voc[iudicium]{iudicium}{judgment} fiat.” Mox venit \voc[tutor]{tutor}{guardian} virginis, vir \voc[iustus]{iustus}{just} et gravis. Cum eo venit etiam \voc[princeps]{princeps}{chief} fori, qui inter mercatores pacem servare solet. Princeps dicit: “Haec \voc[actio]{actio}{action} non parva est. Si furtum probatur, poena erit; sin autem falsa fama est, pudor accusatorem tenebit.”

Tabernarius rem \voc[repeto]{repetit}{to repeat}: “O princeps, ego saepe vidi hunc iuvenem circa mensam meam stare. Panis meus periit, vinum periit, argentum quoque periit. Nemo alius ibi erat; \voc[quisquam]{quisquam}{anyone} eum defendere non potest.” Iuvenis respondet: “Falsum est. Ego nihil abstuli. Si quis verum dicit, audiatur.” Tunc emptor, qui ante amphoras stabat, manum \voc[dexter]{dexteram}{right} tollit et dicit: “Ego hunc iuvenem heri vidi. Non panem rapuit, sed puerum parvum ab igne servavit.”

Princeps iubet omnes tacere. In medio foro ardet \voc[ignis]{ignis}{fire} parvus, nam coquus cibum parat. Emptor dicit: “Heri canis panem rapuit et \voc[usque]{usque}{all the way} ad ignem cucurrit. Iuvenis canem secutus est, panem \voc[excipio]{excepit}{to catch} et puerum, qui prope ignem stabat, servavit. Hoc factum non furtum, sed salus fuit.” Femina addit: “Ego quoque vidi. Iuvenis panem non comedit; panem domino \voc[defero]{detulit}{to bring down/to carry}.” Tabernarius rubet, sed adhuc negat se mentiri.

Tum princeps puerum vocat. Puer parvus, ore \voc[crassus]{crasso}{thick}, sed voce clara, dicit: “Ita est. Canis panem rapuit; iuvenis me iuvit. Si ille non advenisset, ignis mihi nocuisset.” Haec verba omnes movent. Tabernarius, cui \voc[pudor]{pudor}{shame} tandem in pectus descendit, vultum avertit. Princeps dicit: “Nunc res satis patet. Accusator iniustus est, iuvenis autem \voc[victor]{victor}{victor} sine vi est. In hoc \voc[certamen]{certamine}{contest} veritas vicit.”

Sed princeps rem non statim finit. “Restat,” inquit, “quaestio de argento et vino. Emptor vinum emere vult; tabernarius pretium augere vult. Iudicium igitur pretium \voc[decerno]{decernit}{to decide}.” Ille amphoram sumit et vinum inter quinque homines \voc[divido]{dividit}{to divide}. Omnes gustant. Emptor \voc[bibo]{bibit}{to drink} et dicit: “Vinum bonum est, sed pretium nimium.” Tutor quoque bibit et ait: “Pretium \voc[levis]{leve}{light} esse debet, quia vinum non est optimum.” Princeps ridet: “Sic fiat.”

Tabernarius iratus pretium \voc[augeo]{augere}{to increase} conatur, sed princeps eum prohibet. “Non licet,” inquit, “\voc[ullus]{ullam}{any} iniuriam emptori inferre. Si pactum accipis, vinum vende; si non accipis, vendere \voc[desino]{desine}{to stop}.” Tabernarius parumper tacet; deinde dicit: “Pareo.” Sic vinum \voc[vendo]{vendit}{to sell}, emptor autem argentum \voc[emo]{emendo}{to buy} tradit. Quia tabernarius antea pignus postulaverat, emptor annulum parvum ei dederat; nunc princeps iubet annulum reddi. “\voc[pignus]{Pignus}{pledge} iniuste retentum est,” ait.

Tabernarius annulum de mensa \voc[detraho]{detrahit}{to take away} et emptori reddit. Emptor non irascitur, sed dicit: “Si vera verba audis, animus tuus placari potest.” Tabernarius \voc[placo]{placatur}{to appease} et iuveni veniam petit. Iuvenis respondet: “Mihi satis est quod virgo iam non flet.” Virgo subridens ad eum accedit. Tutor dicit: “Si vitam honestam perficis, nuptiae postea \voc[competo]{competent}{to be fitting}.” Iuvenis laetus promittit se laborem non fugere.

Interea femina panem parvum et vinum mixtum parat. Ea pocula \voc[pleo]{plet}{to fill} et hospitibus \voc[tribuo]{tribuit}{to assign} partes aequas. Emptor dicit: “Haec res parva magnum gaudium \voc[efficio]{efficit}{to make}. Ego in foro vinum quaerebam, sed amicitiam inveni.” Virgo stellam ligneam de cera \voc[formo]{format}{to shape}; eam emptori dat et dicit: “Hoc \voc[praemium]{praemium}{reward} accipe, quia nos iuvisti.” Emptor, tali dono laetus, dicit: “Hoc mihi plus valet quam aurum.”

Nox appropinquat. Princeps \voc[totus]{totum}{whole} forum circumspicit et iubet mercatores domum ire, \voc[donec]{donec}{until} silentium fiat. Tabernarius amphoras claudit; emptor vinum portat; iuvenis virginem et matrem eius \voc[comito]{comitatur}{to accompany}. Tutor, qui rem diligenter legerat, tabulas suas claudit; nam multa de lege \voc[lego]{legit}{to read}. In tabulis scribit: “Exceptio falsa non valet; \voc[exceptio]{exceptio}{exception} sine testibus non probatur.” Hoc exemplum postea discipulis narrabit.

In via autem parvus puer iterum occurrit et rogat cur omnes laeti sint. Emptor breviter rem narrat: “Furtum falso dictum est; veritas autem probata est.” Puer exclamat: “O rem pulchram!” Virgo respondet: “Non omnis dies est facilis, sed \voc[tandem]{tandem}{at last} bonus finis venit.” Femina addit: “\voc[finis]{Finis}{end} bonus saepe dolorem sanat.” Iuvenis autem dicit: “Noster animus debet esse fortis.” Omnes consentiunt.

Postero mane tabernarius in foro iterum sedet, sed iam mitior est. Cum emptor ad eum venit, tabernarius dicit: “Amice, heri peccavi. Hodie vinum tibi meliore pretio dabo.” Emptor respondet: “Bene est. Sed memento: qui nimis lucrum quaerit, saepe amicos amittit.” Tabernarius caput movet et dicit: “Vera dicis.” Deinde ad virginem conversus ait: “Tu quoque, puella, noli flere. Hic panis tibi datur.” Virgo panem accipit et matri partem dat.

Sic res in foro Romano composita est. Nemo iam iuvenem hostem putat; immo omnes eum amicum habent. Tabernarius, qui antea severus fuit, nunc ridet; emptor bono vino \voc[fruor]{fruitur}{to enjoy}; femina filiam servat; tutor iudicium laudat; princeps pacem obtinet. Et cum stella vespertina rursus lucet, virgo dicit: “Si homines veritatem audiunt, etiam dies brevis pulcher esse potest.” Ita \voc[posterus]{posteris}{later} diebus omnes hanc fabulam saepe narrant, et memoria eius in foro manet.

\end{readersection}
\end{document}